\chapter{กลศาสตร์คลื่นในบ่อศักย์และปรากฏการณ์ทันเนลลิงเชิงลึก}
\label{chap_potential_wells_tunneling}
\index{บ่อศักย์อนันต์ (Infinite Square Well)}
\index{บ่อศักย์จำกัด (Finite Square Well)}
\index{รอยต่อศักย์ขั้นบันได (Step Potential)}
\index{ควอนตัมทันเนลลิง (Quantum Tunneling)}
\index{การประมาณค่าดับเบิลยูเคบี (WKB Approximation)}
\index{ทฤษฎีการสลายตัวแอลฟาของกามอฟ (Gamow Alpha Decay)}
\index{การสั่นพ้องของการทะลุผ่าน (Transmission Resonance)}
\index{ไดโอดทันเนลลิงแบบสั่นพ้อง (Resonant Tunneling Diode)}
\index{กล้องจุลทรรศน์ส่องกราดแบบอุโมงค์ (STM)}
\index{สเปกโทรสโกปีแบบอุโมงค์ (STS)}
\index{ระเบียบวิธีเมทริกซ์ส่งผ่าน (Transfer Matrix Method)}

% ==========================================================
% แผนบริหารการสอนประจำบทที่ 4
% ==========================================================
\section*{แผนบริหารการสอนประจำบทที่ 4}
\addcontentsline{toc}{section}{แผนบริหารการสอนประจำบทที่ 4}

\begin{lessonplanbox}
\noindent\textbf{หัวข้อเนื้อหาประจำบท}
\begin{enumerate}
    \item รอยต่อศักย์แบบขั้นบันได การสะท้อน การทะลุผ่าน และการอนุรักษ์ฟลักซ์กระแสความน่าจะเป็น
    \item อนุภาคในบ่อศักย์อนันต์หนึ่งมิติ ฐานไอเกนสเตต ความสมบูรณ์ และพลวัตการฟื้นคืนของกลุ่มคลื่น
    \item กล่องศักย์สองมิติและสามมิติ ภาวะเสื่อมของระดับพลังงาน และความหนาแน่นสถานะในโครงสร้างนาโน
    \item บ่อศักย์จำกัดหนึ่งมิติ การวิเคราะห์พาริตี สมการอดิศัย และการแทรกซึมในบริเวณต้องห้ามคลาสสิก
    \item บ่อศักย์และกำแพงศักย์เดลตาของดิแรก สถานะผูกพันเดี่ยว การกระเจิง และแบบจำลองพันธะเคมี
    \item ปรากฏการณ์ควอนตัมทันเนลลิงผ่านกำแพงศักย์สี่เหลี่ยม ปรากฏการณ์แรมเซาเออร์-ทาวน์เซนด์ และไดโอด RTD
    \item การประมาณค่า WKB ทฤษฎีการสลายตัวให้อนุภาคแอลฟาของกามอฟ และปฏิกิริยานิวเคลียร์ฟิวชันในดวงดาว
    \item การประยุกต์ทางเทคโนโลยีล้ำยุค ทฤษฎีเทอร์ซอฟฟ์-ฮามันน์สำหรับกล้อง STM สเปกโทรสโกปี STS และทรานซิสเตอร์ RibbonFET
    \item ห้องปฏิบัติการเสมือนจริง WebAR ปรากฏการณ์ควอนตัมทันเนลลิงและหัวเข็ม STM ระดับอะตอม
    \item การจำลองเชิงตัวเลขด้วยภาษา Python ด้วยระเบียบวิธีเมทริกซ์ส่งผ่าน (Transfer Matrix Method)
    \item ตัวอย่างการแก้โจทย์ปัญหาเชิงวิเคราะห์แบบ P-S-C-T ครอบคลุม 7 สถานการณ์เชิงลึก
    \item สรุปสาระสำคัญประจำบท คำถามเชิงมโนทัศน์ 20 ข้อ และแบบฝึกหัดคำนวณขั้นสูง 20 ข้อ
\end{enumerate}

\noindent\textbf{วัตถุประสงค์เชิงพฤติกรรม}
\begin{enumerate}
    \item คำนวณสัมประสิทธิ์การสะท้อน $R$ และสัมประสิทธิ์การทะลุผ่าน $T$ ที่รอยต่อศักย์ขั้นบันได พร้อมพิสูจน์การอนุรักษ์กระแสความน่าจะเป็น $R + T = 1$ ได้อย่างถูกต้อง
    \item อนุพัทธ์ระดับพลังงานควอนไตซ์ ฟังก์ชันคลื่นไอเกนสเตต ความสัมพันธ์ความไม่แน่นอน และเวลาฟื้นคืนของกลุ่มคลื่นในบ่อศักย์อนันต์ได้อย่างรัดกุม
    \item วิเคราะห์และคำนวณระดับพลังงาน ภาวะเสื่อม และฟังก์ชันความหนาแน่นสถานะ (Density of States) สำหรับโครงสร้างควอนตัม 3D, 2D, 1D และ 0D ได้
    \item แก้สมการอดิศัยของบ่อศักย์จำกัดเพื่อหาระดับพลังงานผูกพันและระยะแทรกซึมเอกซ์โพเนนเชียลได้อย่างแม่นยำ
    \item ประยุกต์ทฤษฎีศักย์เดลตาของดิแรกในการอธิบายสถานะผูกพันและการแยกของระดับพลังงานในระบบสองบ่อศักย์ได้
    \item ประยุกต์การประมาณค่า WKB และทฤษฎีของกามอฟในการคำนวณอัตราการสลายตัวและครึ่งชีวิตของนิวเคลียสกัมมันตรังสีแอลฟาได้อย่างถูกต้อง
    \item อธิบายกลไกทางฟิสิกส์ของกล้อง STM และสเปกโทรสโกปี STS ในการวัดความหนาแน่นสถานะเฉพาะที่ (LDOS) ได้
    \item พัฒนาโปรแกรมภาษา Python ด้วยวิธี Transfer Matrix Method เพื่อคำนวณสเปกตรัมการทะลุผ่านของโครงสร้างกำแพงศักย์ซับซ้อนได้อย่างสมบูรณ์
\end{enumerate}
\end{lessonplanbox}

\newpage

% ==========================================================
% 4.1 รอยต่อศักย์แบบขั้นบันไดและการอนุรักษ์กระแสความน่าจะเป็น
% ==========================================================
\section{รอยต่อศักย์แบบขั้นบันไดและการอนุรักษ์กระแสความน่าจะเป็น}
\label{sec_step_potential}
\index{รอยต่อศักย์ขั้นบันได (Step Potential)}
\index{สัมประสิทธิ์การสะท้อน (Reflection Coefficient)}
\index{สัมประสิทธิ์การทะลุผ่าน (Transmission Coefficient)}

ในการศึกษากลศาสตร์คลื่นควอนตัม ปัญหารอยต่อศักย์แบบขั้นบันได (Step Potential) ถือเป็นระบบพื้นฐานสำคัญที่สุดที่แสดงให้เห็นความแตกต่างอย่างสิ้นเชิงระหว่างกลศาสตร์คลาสสิกของนิวตันกับกลศาสตร์คลื่นควอนตัมของชเรอดิงเงอร์ ฟังก์ชันพลังงานศักย์ของขั้นบันไดหนึ่งมิติที่มีความสูง $V_0$ ณ ตำแหน่ง $x = 0$ เขียนแทนได้ด้วยฟังก์ชันขั้นบันไดเฮวิไซด์ (Heaviside Step Function) ดังนี้
\begin{equation}
V(x) = V_0 \, \Theta(x) = \begin{cases} 
0 & \text{เมื่อ } x < 0 \quad (\text{บริเวณที่ I}) \\ 
V_0 & \text{เมื่อ } x \ge 0 \quad (\text{บริเวณที่ II}) 
\end{cases}
\end{equation}

พิจารณาลำอนุภาคมวล $m$ ที่เคลื่อนที่จากทางซ้าย ($x \to -\infty$) มุ่งหน้าสู่รอยต่อ $x = 0$ ด้วยพลังงานรวม $E$ โดยเราจะแบ่งการวิเคราะห์ออกเป็นสองกรณีหลักตามขนาดของพลังงานเทียบกับความสูงของขั้นบันได

\subsection{กรณีที่หนึ่ง ---- พลังงานสูงกว่าขั้นบันได ($E > V_0$)}
ในกลศาสตร์คลาสสิก เมื่ออนุภาคมีพลังงาน $E > V_0$ เคลื่อนที่ข้ามรอยต่อ $x = 0$ อนุภาคจะเพียงแค่สูญเสียพลังงานจลน์บางส่วนและเคลื่อนที่ต่อไปทางขวาด้วยความเร็วที่ลดลง โดยไม่มีการสะท้อนกลับแม้แต่น้อย ($R_{\text{classical}} = 0$ และ $T_{\text{classical}} = 1$)

ทว่าในกลศาสตร์ควอนตัม สมการชเรอดิงเงอร์ที่ไม่ขึ้นกับเวลาในแต่ละบริเวณมีรูปแบบดังนี้
\begin{align}
\text{บริเวณที่ I } (x < 0) \quad &\colon \quad -\frac{\hbar^2}{2m}\frac{d^2\psi_I}{dx^2} = E\psi_I \quad \Longrightarrow \quad \frac{d^2\psi_I}{dx^2} + k_1^2\psi_I = 0 \\
\text{บริเวณที่ II } (x > 0) \quad &\colon \quad -\frac{\hbar^2}{2m}\frac{d^2\psi_{II}}{dx^2} + V_0\psi_{II} = E\psi_{II} \quad \Longrightarrow \quad \frac{d^2\psi_{II}}{dx^2} + k_2^2\psi_{II} = 0
\end{align}
โดยที่เวกเตอร์คลื่นในแต่ละบริเวณกำหนดโดย
\begin{equation}
k_1 = \frac{\sqrt{2mE}}{\hbar}, \qquad k_2 = \frac{\sqrt{2m(E - V_0)}}{\hbar} < k_1
\end{equation}

ผลเฉลยทั่วไปในบริเวณที่ I ประกอบด้วยคลื่นตกกระทบที่เคลื่อนที่ไปทางขวาและคลื่นสะท้อนที่เคลื่อนที่ไปทางซ้าย ขณะที่ในบริเวณที่ II จะมีเฉพาะคลื่นทะลุผ่านที่เคลื่อนที่ไปทางขวาเท่านั้น (เนื่องจากไม่มีแหล่งกำเนิดคลื่นที่ส่งกลับมาจาก $x \to +\infty$)
\begin{align}
\psi_I(x) &= A e^{ik_1 x} + B e^{-ik_1 x} \\
\psi_{II}(x) &= C e^{ik_2 x}
\end{align}
โดยที่ $A, B, C$ คือแอมพลิจูดเชิงซ้อนของคลื่นตกกระทบ คลื่นสะท้อน และคลื่นทะลุผ่านตามลำดับ

เนื่องจากฟังก์ชันคลื่นและอนุพันธ์อันดับหนึ่งต้องมีความต่อเนื่องที่รอยต่อ $x = 0$ เพื่อรับประกันว่าความหนาแน่นความน่าจะเป็นและกระแสความน่าจะเป็นมีค่าจำกัดและสม่ำเสมอ เราจึงได้เงื่อนไขขอบเขตสองสมการดังนี้
\begin{align}
\psi_I(0) = \psi_{II}(0) \quad &\Longrightarrow \quad A + B = C \\
\left.\frac{d\psi_I}{dx}\right|_{0} = \left.\frac{d\psi_{II}}{dx}\right|_{0} \quad &\Longrightarrow \quad ik_1(A - B) = ik_2 C \quad \Longrightarrow \quad k_1(A - B) = k_2 C
\end{align}

เมื่อแก้ระบบสมการเชิงเส้นเพื่อหาอัตราส่วนของแอมพลิจูดเทียบกับคลื่นตกกระทบ $A$ จะได้
\begin{align}
\frac{B}{A} &= \frac{k_1 - k_2}{k_1 + k_2} \label{eq_step_refl_amp} \\
\frac{C}{A} &= \frac{2k_1}{k_1 + k_2} \label{eq_step_trans_amp}
\end{align}

ในการหาความน่าจะเป็นของการสะท้อนและการทะลุผ่าน เราต้องพิจารณาความหนาแน่นกระแสความน่าจะเป็น (Probability Current Density) ซึ่งนิยามโดย
\begin{equation}
J(x) = \frac{\hbar}{2mi} \left( \psi^* \frac{d\psi}{dx} - \psi \frac{d\psi^*}{dx} \right) = \frac{\hbar}{m} \operatorname{Im}\left(\psi^* \frac{d\psi}{dx}\right)
\end{equation}
เมื่อแทนฟังก์ชันคลื่นของแต่ละองค์ประกอบลงไป จะได้ฟลักซ์กระแสความน่าจะเป็นดังนี้
\begin{equation}
J_{\text{inc}} = \frac{\hbar k_1}{m} |A|^2, \qquad J_{\text{ref}} = \frac{\hbar k_1}{m} |B|^2, \qquad J_{\text{trans}} = \frac{\hbar k_2}{m} |C|^2
\end{equation}

ดังนั้น สัมประสิทธิ์การสะท้อน (Reflection Coefficient, $R$) และสัมประสิทธิ์การทะลุผ่าน (Transmission Coefficient, $T$) จึงนิยามจากอัตราส่วนของฟลักซ์กระแสความน่าจะเป็น
\begin{align}
R &\equiv \frac{J_{\text{ref}}}{J_{\text{inc}}} = \left| \frac{B}{A} \right|^2 = \left( \frac{k_1 - k_2}{k_1 + k_2} \right)^2 \label{eq_step_refl_coeff} \\
T &\equiv \frac{J_{\text{trans}}}{J_{\text{inc}}} = \frac{k_2}{k_1} \left| \frac{C}{A} \right|^2 = \frac{k_2}{k_1} \left( \frac{2k_1}{k_1 + k_2} \right)^2 = \frac{4k_1 k_2}{(k_1 + k_2)^2} \label{eq_step_trans_coeff}
\end{align}

เมื่อนำสัมประสิทธิ์ทั้งสองมารวมกัน จะได้ผลลัพธ์ที่สอดคล้องกับกฎทรงความน่าจะเป็นอย่างสมบูรณ์
\begin{equation}
R + T = \frac{(k_1 - k_2)^2 + 4k_1 k_2}{(k_1 + k_2)^2} = \frac{k_1^2 - 2k_1 k_2 + k_2^2 + 4k_1 k_2}{(k_1 + k_2)^2} = \frac{(k_1 + k_2)^2}{(k_1 + k_2)^2} = 1
\end{equation}

สมการข้างต้นแสดงให้เห็นอย่างชัดเจนว่า แม้พลังงานของอนุภาคจะสูงกว่าความสูงของกำแพงศักย์ ($E > V_0$) อนุภาคในกลศาสตร์ควอนตัมก็ยังมีโอกาสที่จะ \textbf{สะท้อนกลับ} จากรอยต่อเสมอ ตราบใดที่พลังงานศักย์เกิดการเปลี่ยนแปลงอย่างฉับพลัน ปรากฏการณ์นี้สอดคล้องโดยตรงกับการสะท้อนของคลื่นแม่เหล็กไฟฟ้าเมื่อเดินทางผ่านรอยต่อระหว่างตัวกลางที่มีดัชนีหักเหแสงต่างกันตามสมการเฟรสเนล (Fresnel Equations)

\subsection{กรณีที่สอง ---- พลังงานต่ำกว่าขั้นบันได ($E < V_0$) และคลื่นสลายตัว}
ในกรณีที่อนุภาคมีพลังงานน้อยกว่าความสูงของขั้นบันได ($E < V_0$) กลศาสตร์คลาสสิกจะทำนายว่าบริเวณ $x > 0$ เป็นบริเวณต้องห้าม (Classically Forbidden Region) เนื่องจากพลังงานจลน์ $K = E - V_0 < 0$ ซึ่งเป็นไปไม่ได้ในทางกายภาพ ทำให้อนุภาคต้องสะท้อนกลับ $100\%$ ที่ตำแหน่ง $x = 0$

สำหรับกลศาสตร์ควอนตัม ในบริเวณที่ II สมการชเรอดิงเงอร์จะกลายเป็น
\begin{equation}
\frac{d^2\psi_{II}}{dx^2} - \kappa^2 \psi_{II} = 0, \quad \text{โดยที่ } \kappa = \frac{\sqrt{2m(V_0 - E)}}{\hbar} \in \mathbb{R}^+
\end{equation}
ผลเฉลยทั่วไปคือ $\psi_{II}(x) = C e^{-\kappa x} + D e^{\kappa x}$ แต่เนื่องจากฟังก์ชันคลื่นต้องมีขอบเขตจำกัดทางกายภาพและสามารถนอร์มอลไลซ์ได้ เมื่อ $x \to +\infty$ เทอม $e^{\kappa x}$ จะลู่ออกสู่ค่าอนันต์ ดังนั้นค่าคงที่ $D$ จึงต้องเป็นศูนย์อย่างเคร่งครัด ทำให้ฟังก์ชันคลื่นในบริเวณที่สองลดลงแบบเอกซ์โพเนนเชียล (Evanescent Wave)
\begin{equation}
\psi_{II}(x) = C e^{-\kappa x} \quad (x \ge 0)
\end{equation}

เมื่อนำเงื่อนไขความต่อเนื่องที่ $x = 0$ มาประยุกต์ใช้
\begin{align}
\psi_I(0) = \psi_{II}(0) \quad &\Longrightarrow \quad A + B = C \\
\left.\frac{d\psi_I}{dx}\right|_{0} = \left.\frac{d\psi_{II}}{dx}\right|_{0} \quad &\Longrightarrow \quad ik_1(A - B) = -\kappa C
\end{align}
แก้สมการหาแอมพลิจูดการสะท้อน $B/A$
\begin{equation}
\frac{B}{A} = \frac{k_1 - i\kappa}{k_1 + i\kappa} = \frac{k_1 - i\kappa}{\sqrt{k_1^2 + \kappa^2}} \Big/ \frac{k_1 + i\kappa}{\sqrt{k_1^2 + \kappa^2}} = e^{-2i\phi}
\end{equation}
โดยที่เฟสชิฟต์ $\phi$ สอดคล้องกับ $\tan\phi = \frac{\kappa}{k_1} = \sqrt{\frac{V_0 - E}{E}}$

สัมประสิทธิ์การสะท้อนในกรณีนี้คือ
\begin{equation}
R = \left| \frac{B}{A} \right|^2 = |e^{-2i\phi}|^2 = 1
\end{equation}
นั่นคือ อนุภาคจะสะท้อนกลับอย่างสมบูรณ์แบบ $100\%$ เช่นเดียวกับการทำนายทางคลาสสิก ทว่าความแตกต่างอันลึกซึ้งอยู่ที่ \textbf{ฟังก์ชันคลื่นไม่ได้เป็นศูนย์ทันทีที่รอยต่อ} แต่สามารถแทรกซึมเข้าไปในเนื้อกำแพงศักย์ได้ลึกในระดับหนึ่ง โดยความหนาแน่นความน่าจะเป็นที่จะพบอนุภาคในบริเวณต้องห้ามคลาสสิก ($x > 0$) มีค่าเท่ากับ
\begin{equation}
|\psi_{II}(x)|^2 = |C|^2 e^{-2\kappa x} = \frac{4k_1^2}{k_1^2 + \kappa^2} |A|^2 e^{-2\kappa x}
\end{equation}

เรานิยาม \textbf{ระยะแทรกซึมลักษณะเฉพาะ (Characteristic Penetration Depth, $\delta_p$)} ว่าเป็นระยะทางที่ความหนาแน่นความน่าจะเป็นลดลงเหลือ $1/e$ ของค่าที่ผิวหน้า
\begin{equation}
\delta_p = \frac{1}{2\kappa} = \frac{\hbar}{2\sqrt{2m(V_0 - E)}}
\end{equation}
เมื่อคำนวณกระแสความน่าจะเป็นในบริเวณที่ II จะพบว่า
\begin{equation}
J_{II} = \frac{\hbar}{m} \operatorname{Im}\left( \psi_{II}^* \frac{d\psi_{II}}{dx} \right) = \frac{\hbar}{m} \operatorname{Im}\left( C^* e^{-\kappa x} (-\kappa C e^{-\kappa x}) \right) = -\frac{\hbar\kappa}{m} |C|^2 e^{-2\kappa x} \operatorname{Im}(1) = 0
\end{equation}
การที่ $J_{II} = 0$ ยืนยันว่าไม่มีการไหลถ่ายโอนฟลักซ์ความน่าจะเป็นสุทธิเข้าไปในเนื้อกำแพงอย่างถาวร แต่เป็นเพียงการสถิตของคลื่นระเหยหาย (Evanescent Mode) ซึ่งคลื่นที่แทรกซึมเข้าไปจะไหลวนกลับออกมาสู่บริเวณที่ I ทั้งหมด ก่อให้เกิดปรากฏการณ์การเลื่อนตามขวางกูส์-แฮนเชน (Goos-Hänchen Shift) ของกลุ่มคลื่นสสาร เช่นเดียวกับที่เกิดขึ้นในการสะท้อนกลับหมดของแสงเลเซอร์

% ==========================================================
% 4.2 อนุภาคในบ่อศักย์อนันต์หนึ่งมิติและโครงสร้างพีชคณิตไอเกนสเตต
% ==========================================================
\section{อนุภาคในบ่อศักย์อนันต์หนึ่งมิติและโครงสร้างพีชคณิตไอเกนสเตต}
\label{sec_infinite_well}
\index{บ่อศักย์อนันต์หนึ่งมิติ (1D Infinite Well)}
\index{ระดับพลังงานควอนไตซ์ (Quantized Energy Levels)}
\index{การฟื้นคืนของกลุ่มคลื่น (Quantum Revivals)}

แบบจำลองอนุภาคในกล่องศักย์อนันต์ (Particle in a Box หรือ Infinite Square Well) ถือเป็นระบบควอนตัมขั้นปฐมภูมิที่สาธิตหลักการควอนไตเซชันของพลังงานอันเป็นผลพวงโดยตรงจากการกักขังเชิงพื้นที่ (Spatial Confinement)

\subsection{การแก้สมการและการควอนไตซ์ของพลังงาน}
กำหนดฟังก์ชันพลังงานศักย์หนึ่งมิติตามแนวแกน $x$ ดังนี้
\begin{equation}
V(x) = \begin{cases}
0 & \text{เมื่อ } 0 < x < L \\
\infty & \text{เมื่อ } x \le 0 \text{ หรือ } x \ge L
\end{cases}
\end{equation}
เนื่องจากกำแพงศักย์ภายนอกบ่อมีค่าเป็นอนันต์ พลังงานที่ต้องใช้ในการผลักให้อนุภาคหลุดออกไปนอกบ่อย่อมเป็นอนันต์ ความน่าจะเป็นในการตรวจพบอนุภาคภายนอกบ่อจึงต้องเป็นศูนย์อย่างสมบูรณ์ นั่นคือ
\begin{equation}
\psi(x) = 0 \quad \text{สำหรับทุกค่า } x \le 0 \text{ และ } x \ge L
\end{equation}

ภายในบ่อศักย์ ($0 < x < L$) พลังงานศักย์ $V(x) = 0$ สมการชเรอดิงเงอร์ที่ไม่ขึ้นกับเวลาลดรูปลงเป็น
\begin{equation}
-\frac{\hbar^2}{2m}\frac{d^2\psi(x)}{dx^2} = E\psi(x) \quad \Longrightarrow \quad \frac{d^2\psi(x)}{dx^2} + k^2\psi(x) = 0
\end{equation}
โดยที่ $k = \frac{\sqrt{2mE}}{\hbar} > 0$ ผลเฉลยทั่วไปคือ
\begin{equation}
\psi(x) = A \sin(kx) + B \cos(kx)
\end{equation}

การบังคับใช้เงื่อนไขความต่อเนื่องของฟังก์ชันคลื่นที่ขอบเขตทั้งสองด้านนำไปสู่ผลลัพธ์ดังนี้
\begin{enumerate}
    \item ที่ขอบด้านซ้าย $x = 0$ ต้องสอดคล้องกับ $\psi(0) = 0$ จะได้
    \begin{equation}
    A \sin(0) + B \cos(0) = 0 \quad \Longrightarrow \quad B = 0
    \end{equation}
    ทำให้ฟังก์ชันคลื่นลดรูปเหลือ $\psi(x) = A \sin(kx)$
    
    \item ที่ขอบด้านขวา $x = L$ ต้องสอดคล้องกับ $\psi(L) = 0$ จะได้
    \begin{equation}
    A \sin(kL) = 0
    \end{equation}
    เนื่องจาก $A \ne 0$ (มิฉะนั้นฟังก์ชันคลื่นจะเป็นศูนย์ทุกหนแห่ง ซึ่งขัดแย้งกับการมีอยู่ของอนุภาค) ดังนั้นจึงบังคับให้อาร์กิวเมนต์ของฟังก์ชันไซน์ต้องเป็นจำนวนเต็มเท่าของ $\pi$ นั่นคือ
    \begin{equation}
    k_n L = n\pi \quad \Longrightarrow \quad k_n = \frac{n\pi}{L}, \qquad n = 1, 2, 3, \dots
    \end{equation}
\end{enumerate}

\begin{figure}[H]
\centering
\begin{tikzpicture}[scale=1.15]
    % Potential Well Walls
    \fill[slateGrey!30!white] (-0.8, 0) rectangle (0, 4.5);
    \fill[slateGrey!30!white] (5, 0) rectangle (5.8, 4.5);
    \draw[darkNavy, line width=2.5pt] (0, 4.5) -- (0, 0) -- (5, 0) -- (5, 4.5);
    \node at (-0.4, 2.5) [rotate=90, font=\bfseries\color{darkNavy}] {$V(x) = \infty$};
    \node at (5.4, 2.5) [rotate=-90, font=\bfseries\color{darkNavy}] {$V(x) = \infty$};
    
    % Ground Level n=1
    \draw[dashed, slateGrey!50] (0, 0.8) -- (5, 0.8) node[right, font=\footnotesize] {$E_1$};
    \draw[cyberBlue, line width=1.8pt, domain=0:5, samples=80] plot (\x, {0.8 + 0.5*sin(180*\x/5)});
    \node at (2.5, 1.5) [font=\small\bfseries\color{cyberBlue}] {$n=1, \psi_1(x)$};
    
    % Level n=2
    \draw[dashed, slateGrey!50] (0, 2.0) -- (5, 2.0) node[right, font=\footnotesize] {$E_2 = 4E_1$};
    \draw[quantumGold, line width=1.8pt, domain=0:5, samples=80] plot (\x, {2.0 + 0.5*sin(360*\x/5)});
    \node at (1.25, 2.65) [font=\small\bfseries\color{quantumGold}] {$n=2$};
    
    % Level n=3
    \draw[dashed, slateGrey!50] (0, 3.6) -- (5, 3.6) node[right, font=\footnotesize] {$E_3 = 9E_1$};
    \draw[rubyRed, line width=1.8pt, domain=0:5, samples=100] plot (\x, {3.6 + 0.5*sin(540*\x/5)});
    \node at (0.8, 4.25) [font=\small\bfseries\color{rubyRed}] {$n=3$};
    
    % Axes and labels
    \draw[-{Stealth}, line width=1.2pt] (0, -0.4) -- (5.5, -0.4) node[right] {$x$};
    \node at (0, -0.7) [font=\small\bfseries] {$0$};
    \node at (5, -0.7) [font=\small\bfseries] {$L$};
\end{tikzpicture}
\caption{ฟังก์ชันคลื่นสถานะเจาะจง $\psi_n(x)$ สามระดับแรกในบ่อศักย์อนันต์หนึ่งมิติ}
\label{fig_infinite_well_wavefunctions}
\end{figure}

ข้อสังเกตที่สำคัญยิ่งคือ $n$ ไม่สามารถเป็นศูนย์ได้ ($n \ne 0$) เพราะหาก $n = 0$ จะส่งผลให้ $k_0 = 0$ และ $\psi_0(x) = 0$ ทุกค่า $x$ ซึ่งหมายความว่าไม่มีอนุภาคอยู่ในระบบ และค่า $n$ ที่เป็นลบจะให้เพียงฟังก์ชันคลื่นเดิมที่คูณด้วยเครื่องหมายลบ ซึ่งเป็นสถานะทางกายภาพเดียวกัน ดังนั้น เลขควอนตัม $n$ จึงต้องเป็นจำนวนเต็มบวกเท่านั้น ($n \in \mathbb{N}$)

จากความสัมพันธ์ $E = \frac{\hbar^2 k^2}{2m}$ ระดับพลังงานที่ยอมรับได้ทางควอนตัมจึงถูกควอนไตซ์อย่างเคร่งครัด
\begin{equation}
E_n = \frac{\hbar^2 k_n^2}{2m} = \frac{n^2 \pi^2 \hbar^2}{2mL^2} = n^2 E_1, \qquad n = 1, 2, 3, \dots
\end{equation}
โดยที่พลังงานสถานะพื้น (Ground State Energy, $n=1$) มีค่าเท่ากับ
\begin{equation}
E_1 = \frac{\pi^2 \hbar^2}{2mL^2} > 0
\end{equation}
การที่พลังงานสถานะพื้นมีค่ามากกว่าศูนย์อย่างเด็ดขาด เรียกว่า \textbf{พลังงานจุดศูนย์ (Zero-Point Energy)} ซึ่งสอดคล้องอย่างสมบูรณ์กับหลักความไม่แน่นอนของไฮเซนแบร์ก เนื่องจากอนุภาคถูกจำกัดตำแหน่งอยู่ในช่วงความกว้าง $L$ ความไม่แน่นอนของตำแหน่งจึงมีขอบเขตจำกัด $\Delta x \le L$ ส่งผลให้โมเมนตัมต้องมีความไม่แน่นอนอย่างน้อยที่สุด $\Delta p \ge \frac{\hbar}{2L}$ และทำให้พลังงานจลน์ไม่สามารถเป็นศูนย์ได้

เมื่อคำนวณค่านอร์มอลไลเซชัน $\int_0^L |\psi_n(x)|^2 dx = 1$
\begin{equation}
|A|^2 \int_0^L \sin^2\left(\frac{n\pi x}{L}\right) dx = |A|^2 \frac{L}{2} = 1 \quad \Longrightarrow \quad A = \sqrt{\frac{2}{L}}
\end{equation}
จะได้ฟังก์ชันคลื่นไอเกนสเตตที่สมบูรณ์คือ
\begin{equation}
\psi_n(x) = \sqrt{\frac{2}{L}} \sin\left(\frac{n\pi x}{L}\right), \qquad n = 1, 2, 3, \dots
\end{equation}

\subsection{ความสัมพันธ์การตั้งฉาก ความสมบูรณ์ และการคำนวณค่าคาดหวัง}
เซตของฟังก์ชันคลื่นไอเกนสเตต $\{\psi_n(x)\}$ ก่อร่างเป็นฐานตั้งฉากปรกติสมบูรณ์ (Complete Orthonormal Basis) ในปริภูมิฮิลเบิร์ต $L^2(0, L)$ ซึ่งสอดคล้องกับคุณสมบัติดังนี้
\begin{align}
\text{ความตั้งฉากปรกติ} \quad &\colon \quad \int_0^L \psi_m^*(x)\psi_n(x) dx = \delta_{mn} \\
\text{ความสมบูรณ์} \quad &\colon \quad \sum_{n=1}^\infty \psi_n(x)\psi_n^*(x') = \delta(x - x')
\end{align}

ดังนั้น ฟังก์ชันสถานะคลื่นใดๆ $\Psi(x, 0)$ ในบ่อศักย์สามารถกระจายตัวในรูปอนุกรมฟูเรียร์ของไอเกนสเตตได้เสมอ
\begin{equation}
\Psi(x, 0) = \sum_{n=1}^\infty c_n \psi_n(x), \quad \text{โดยที่ } c_n = \int_0^L \psi_n^*(x)\Psi(x, 0) dx
\end{equation}
และวิวัฒนาการตามเวลาของสถานะควอนตัมจะถูกกำหนดโดย
\begin{equation}
\Psi(x, t) = \sum_{n=1}^\infty c_n \psi_n(x) e^{-iE_n t/\hbar} = \sqrt{\frac{2}{L}} \sum_{n=1}^\infty c_n \sin\left(\frac{n\pi x}{L}\right) \exp\left( -i \frac{n^2 \pi^2 \hbar t}{2mL^2} \right)
\end{equation}

เมื่อคำนวณค่าคาดหวังของปริมาณทางกายภาพในสถานะไอเกนสเตต $\psi_n(x)$ จะได้ผลลัพธ์เชิงวิเคราะห์ที่รัดกุมดังนี้
\begin{align}
\langle x \rangle_n &= \int_0^L x |\psi_n(x)|^2 dx = \frac{2}{L}\int_0^L x \sin^2\left(\frac{n\pi x}{L}\right) dx = \frac{L}{2} \\
\langle x^2 \rangle_n &= \frac{2}{L}\int_0^L x^2 \sin^2\left(\frac{n\pi x}{L}\right) dx = L^2 \left( \frac{1}{3} - \frac{1}{2n^2\pi^2} \right) \\
\langle p_x \rangle_n &= \int_0^L \psi_n^*(x) \left(-i\hbar\frac{d}{dx}\right) \psi_n(x) dx = 0 \\
\langle p_x^2 \rangle_n &= 2m E_n = \frac{n^2\pi^2\hbar^2}{L^2}
\end{align}

ความไม่แน่นอนของตำแหน่งและโมเมนตัมจึงคำนวณได้เป็น
\begin{align}
\Delta x &= \sqrt{\langle x^2 \rangle - \langle x \rangle^2} = L \sqrt{\frac{1}{12} - \frac{1}{2n^2\pi^2}} \\
\Delta p &= \sqrt{\langle p^2 \rangle - \langle p \rangle^2} = \frac{n\pi\hbar}{L}
\end{align}
และผลคูณความไม่แน่นอนมีค่าเท่ากับ
\begin{equation}
\Delta x \Delta p = \hbar \sqrt{\frac{n^2\pi^2}{12} - \frac{1}{2}} = \frac{\hbar}{2} \sqrt{\frac{n^2\pi^2}{3} - 2}
\end{equation}
สำหรับสถานะพื้น ($n=1$) จะได้ $\Delta x \Delta p = \frac{\hbar}{2} \sqrt{\frac{\pi^2}{3} - 2} \approx 0.568 \hbar > \frac{\hbar}{2}$ ซึ่งสอดคล้องกับหลักความไม่แน่นอนของไฮเซนแบร์กอย่างเคร่งครัด

\subsection{ปรากฏการณ์การฟื้นคืนของกลุ่มคลื่นควอนตัม (Quantum Revivals)}
พิจารณาวิวัฒนาการตามเวลาของกลุ่มคลื่น $\Psi(x, t)$ จะเห็นได้ว่าเฟสของแต่ละองค์ประกอบขึ้นอยู่กับ $n^2$ ผ่านเทอม $\exp(-i\omega_n t)$ โดยที่ $\omega_n = \frac{E_n}{\hbar} = n^2 \omega_1$ เมื่อ $\omega_1 = \frac{\pi^2\hbar}{2mL^2}$

เนื่องจากความถี่ของทุกฮาร์มอนิกเป็นสัดส่วนกับกำลังสองของจำนวนเต็ม ($n^2$) เมื่อเวลาผ่านไปจนกระทั่งทุกเฟสหมุนครบรอบ $2\pi$ พร้อมกัน กลุ่มคลื่นที่เคยแผ่กระจายตัวออกจะกลับมารวมตัวกันเป็นรูปร่างตั้งต้นอย่างสมบูรณ์แบบ เราเรียกเวลานี้ว่า \textbf{เวลาฟื้นคืนควอนตัม (Quantum Revival Time, $T_{\text{rev}}$)}
\begin{equation}
\omega_1 T_{\text{rev}} = 2\pi \quad \Longrightarrow \quad T_{\text{rev}} = \frac{2\pi}{\omega_1} = \frac{4mL^2}{\pi\hbar}
\end{equation}
ณ เวลา $t = T_{\text{rev}}$ ฟังก์ชันคลื่นจะคืนสภาพเดิม $\Psi(x, T_{\text{rev}}) = \Psi(x, 0)$ ยิ่งไปกว่านั้น ณ เวลาเศษส่วน เช่น $t = T_{\text{rev}}/2$ กลุ่มคลื่นจะเกิดการฟื้นคืนแบบกระจกเงา $\Psi(x, T_{\text{rev}}/2) = -\Psi(L-x, 0)$ และที่ $t = T_{\text{rev}}/4$ กลุ่มคลื่นจะแยกออกเป็นสองร่างคู่ขนานที่พันกันทางควอนตัม เรียกว่าสถานะแมวชเรอดิงเงอร์ (Schrödinger Cat State) ปรากฏการณ์นี้ได้รับการยืนยันการทดลองจริงในอะตอมริดเบิร์ก (Rydberg Atoms) และการทดลองทัลบอตของคลื่นสสารในคริสต์ศตวรรษที่ 21

% ==========================================================
% 4.3 กล่องศักย์สองมิติและสามมิติ: ภาวะเสื่อมและความหนาแน่นสถานะ
% ==========================================================
\section{กล่องศักย์สองมิติและสามมิติ ภาวะเสื่อม และความหนาแน่นสถานะ}
\label{sec_multidimensional_wells}
\index{กล่องศักย์สามมิติ (3D Quantum Box)}
\index{ภาวะเสื่อม (Degeneracy)}
\index{ความหนาแน่นสถานะ (Density of States)}
\index{โครงสร้างบ่อควอนตัม (Quantum Well 2DEG)}

ในการประยุกต์ใช้วิชาฟิสิกส์ควอนตัมกับฟิสิกส์สถานะของแข็งและอุปกรณ์สารกึ่งตัวนำ อนุภาคจะเคลื่อนที่ในปริภูมิหลายมิติ ซึ่งนำไปสู่ปรากฏการณ์สำคัญสองประการ ได้แก่ ภาวะเสื่อมของระดับพลังงาน (Energy Degeneracy) และพฤติกรรมของฟังก์ชันความหนาแน่นสถานะ (Density of States) ในมิติต่างๆ

\subsection{กล่องศักย์สี่เหลี่ยมผืนผ้าสองมิติและภาวะเสื่อมเชิงสมมาตร}
พิจารณาอนุภาคมวล $m$ ในกล่องศักย์สองมิติที่มีความกว้างตามแนวแกน $x$ เท่ากับ $L_x$ และตามแนวแกน $y$ เท่ากับ $L_y$ สมการชเรอดิงเงอร์ที่ไม่ขึ้นกับเวลาคือ
\begin{equation}
-\frac{\hbar^2}{2m}\left(\frac{\partial^2\psi}{\partial x^2} + \frac{\partial^2\psi}{\partial y^2}\right) = E\psi(x, y)
\end{equation}
โดยใช้ระเบียบวิธีแยกตัวแปร (Separation of Variables) $\psi(x, y) = X(x)Y(y)$ จะได้ผลเฉลยไอเกนสเตตและระดับพลังงาน
\begin{align}
\psi_{n_x, n_y}(x, y) &= \sqrt{\frac{4}{L_x L_y}} \sin\left(\frac{n_x\pi x}{L_x}\right) \sin\left(\frac{n_y\pi y}{L_y}\right) \\
E_{n_x, n_y} &= \frac{\pi^2\hbar^2}{2m}\left(\frac{n_x^2}{L_x^2} + \frac{n_y^2}{L_y^2}\right), \qquad n_x, n_y \in \{1, 2, 3, \dots\}
\end{align}

ในกรณีพิเศษที่กล่องมีสมมาตรแบบจัตุรัส ($L_x = L_y = L$) ระดับพลังงานจะลดรูปเหลือ
\begin{equation}
E_{n_x, n_y} = \frac{\pi^2\hbar^2}{2mL^2}(n_x^2 + n_y^2)
\end{equation}
เมื่อพิจารณาสถานะ $(n_x, n_y) = (1, 2)$ และ $(2, 1)$ จะพบว่า
\begin{equation}
E_{1, 2} = E_{2, 1} = \frac{\pi^2\hbar^2}{2mL^2}(1^2 + 2^2) = 5 E_1
\end{equation}
ทว่าฟังก์ชันคลื่นทั้งสองกลับมีรูปร่างที่แตกต่างกันอย่างสิ้นเชิง
\begin{equation}
\psi_{1, 2}(x, y) = \frac{2}{L}\sin\left(\frac{\pi x}{L}\right)\sin\left(\frac{2\pi y}{L}\right) \quad \ne \quad \psi_{2, 1}(x, y) = \frac{2}{L}\sin\left(\frac{2\pi x}{L}\right)\sin\left(\frac{\pi y}{L}\right)
\end{equation}
สภาวะที่ไอเกนสเตตที่แตกต่างกันทางกายภาพตั้งแต่สองสถานะขึ้นไปมีค่าพลังงานเจาะจงเท่ากันทุกประการนี้ เรียกว่า \textbf{ภาวะเสื่อม (Degeneracy)} โดยภาวะเสื่อมระหว่าง $(1, 2)$ และ $(2, 1)$ เกิดจากสมมาตรการสลับแกนพิกัด (Geometric Symmetry) ขณะที่ภาวะเสื่อมเช่น $(1, 7)$ และ $(5, 5)$ ซึ่งให้ผลรวมกำลังสองเท่ากับ $50$ ($1^2 + 7^2 = 5^2 + 5^2 = 50$) เกิดจากคุณสมบัติทางทฤษฎีจำนวน เรียกว่า \textbf{ภาวะเสื่อมโดยบังเอิญ (Accidental Degeneracy)}

\subsection{ความหนาแน่นสถานะ (Density of States) ในโครงสร้างนาโนควอนตัม}
ฟังก์ชันความหนาแน่นสถานะ $\rho(E) = \frac{dN}{dE}$ นิยามเป็นจำนวนสถานะควอนตัมต่อหน่วยช่วงพลังงาน ซึ่งเป็นตัวแปรชี้ขาดพฤติกรรมการนำไฟฟ้า การเปล่งแสงเลเซอร์ และสมบัติทางเทอร์โมไดนามิกส์ของสสาร

\begin{table}[H]
\centering
\caption{การเปรียบเทียบฟังก์ชันความหนาแน่นสถานะในมิติต่างๆ ของโครงสร้างควอนตัม}
\label{tab_density_of_states}
\begin{tabular}{llll}
\toprule
\textbf{มิติของระบบ} & \textbf{โครงสร้างตัวอย่าง} & \textbf{ฟังก์ชันความหนาแน่นสถานะ $\rho(E)$} & \textbf{ลักษณะทางเรขาคณิต} \\
\midrule
3 มิติ (Bulk 3D) & ผลึกสารกึ่งตัวนำก้อนใหญ่ & $\rho_{3D}(E) = \frac{(2m^*)^{3/2}}{2\pi^2\hbar^3}\sqrt{E}$ & พาราโบลา ($\propto \sqrt{E}$) \\
2 มิติ (Quantum Well) & ก๊าซอิเล็กตรอน 2D (GaAs/AlGaAs) & $\rho_{2D}(E) = \frac{m^*}{\pi\hbar^2}\sum_n \Theta(E - E_n)$ & ขั้นบันไดคงที่ (Staircase) \\
1 มิติ (Quantum Wire) & ท่อนาโนคาร์บอน เส้นลวดนาโน & $\rho_{1D}(E) = \frac{\sqrt{2m^*}}{\pi\hbar}\sum_n \frac{\Theta(E - E_n)}{\sqrt{E - E_n}}$ & พุ่งชันสู่ค่าอนันต์ (van Hove) \\
0 มิติ (Quantum Dot) & จุดควอนตัม (อะตอมประดิษฐ์) & $\rho_{0D}(E) = 2\sum_i \delta(E - E_i)$ & พีคเดลตาไม่ต่อเนื่อง ($\delta$) \\
\bottomrule
\end{tabular}
\end{table}

ในอุปกรณ์เลเซอร์สารกึ่งตัวนำยุคใหม่ การลดมิติของระบบจาก 3D ลงสู่บ่อควอนตัม 2D และจุดควอนตัม 0D ช่วยปรับเปลี่ยนความหนาแน่นสถานะให้มีความหนาแน่นสูงเป็นพิเศษ ณ ระดับพลังงานที่ต้องการ ส่งผลให้กระแสไฟฟ้าเริ่มต้นในการทำงานของเลเซอร์ (Threshold Current) ลดต่ำลงอย่างมหาศาล และมีประสิทธิภาพควอนตัมสูงขึ้นอย่างเด่นชัด

% ==========================================================
% 4.4 บ่อศักย์จำกัดหนึ่งมิติและสมการอดิศัย
% ==========================================================
\section{บ่อศักย์จำกัดหนึ่งมิติและสมการอดิศัย}
\label{sec_finite_well}
\index{บ่อศักย์จำกัด (Finite Square Well)}
\index{สมการอดิศัย (Transcendental Equation)}
\index{พาริตี (Parity)}

ในความเป็นจริงทางกายภาพ กำแพงศักย์ไม่มีทางที่จะสูงเป็นอนันต์ได้ อนุภาคจะถูกกักขังอยู่ในบ่อศักย์จำกัด (Finite Square Well) ซึ่งมีฟังก์ชันพลังงานศักย์สมมาตรดังนี้
\begin{equation}
V(x) = \begin{cases}
-V_0 & \text{เมื่อ } |x| \le a \quad (\text{ภายในบ่อ}) \\
0 & \text{เมื่อ } |x| > a \quad (\text{ภายนอกบ่อ})
\end{cases}
\end{equation}
โดยที่ $V_0 > 0$ คือความลึกของบ่อศักย์ และความกว้างของบ่อคือ $2a$

เราสนใจสภาวะสถานะผูกพัน (Bound States) ซึ่งอนุภาคถูกกักขังอยู่ในบ่อ นั่นคือมีพลังงานอยู่ในช่วง $-V_0 < E < 0$ นิยามเวกเตอร์คลื่นภายในและภายนอกบ่อ
\begin{align}
k &= \frac{\sqrt{2m(E + V_0)}}{\hbar} > 0 \quad (\text{ภายในบ่อ } |x| \le a) \\
\kappa &= \frac{\sqrt{-2mE}}{\hbar} = \frac{\sqrt{2m|E|}}{\hbar} > 0 \quad (\text{ภายนอกบ่อ } |x| > a)
\end{align}

เนื่องจากฮามิลโทเนียนมีความสมมาตรต่อการสะท้อนพิกัด $V(-x) = V(x)$ ตัวดำเนินการพาริตี $\hat{\Pi}$ จะสลับที่กับฮามิลโทเนียน $[\hat{H}, \hat{\Pi}] = 0$ ฟังก์ชันคลื่นไอเกนสเตตจึงสามารถแยกออกเป็นสถานะพาริตีคู่ (Even Parity) และสถานะพาริตีคี่ (Odd Parity) ได้อย่างอิสระ

\subsection{สถานะพาริตีคู่ (Even Parity States)}
สำหรับสถานะสมมาตรคู่ $\psi(-x) = \psi(x)$ ฟังก์ชันคลื่นที่สอดคล้องกับเงื่อนไขการลู่เข้าสู่ศูนย์ที่ระยะอนันต์คือ
\begin{equation}
\psi_{\text{even}}(x) = \begin{cases}
C e^{\kappa x} & \text{เมื่อ } x < -a \\
A \cos(kx) & \text{เมื่อ } -a \le x \le a \\
C e^{-\kappa x} & \text{เมื่อ } x > a
\end{cases}
\end{equation}

นำเงื่อนไขความต่อเนื่องของฟังก์ชันคลื่นและอนุพันธ์อันดับหนึ่งที่ตำแหน่งรอยต่อ $x = a$ มาประยุกต์
\begin{align}
A \cos(ka) &= C e^{-\kappa a} \\
-k A \sin(ka) &= -\kappa C e^{-\kappa a}
\end{align}
เมื่อนำสมการทั้งสองมาหารกัน จะได้สมการเงื่อนไขพลังงานสำหรับสถานะพาริตีคู่
\begin{equation}
ka \tan(ka) = \kappa a
\end{equation}

\subsection{สถานะพาริตีคี่ (Odd Parity States)}
สำหรับสถานะไม่สมมาตร $\psi(-x) = -\psi(x)$ ฟังก์ชันคลื่นกำหนดโดย
\begin{equation}
\psi_{\text{odd}}(x) = \begin{cases}
-D e^{\kappa x} & \text{เมื่อ } x < -a \\
B \sin(kx) & \text{เมื่อ } -a \le x \le a \\
D e^{-\kappa x} & \text{เมื่อ } x > a
\end{cases}
\end{equation}
เมื่อจับคู่ความต่อเนื่องที่ $x = a$ และหารสมการเข้าด้วยกัน จะได้สมการเงื่อนไขพลังงานสำหรับสถานะพาริตีคี่
\begin{equation}
-ka \cot(ka) = \kappa a
\end{equation}

\subsection{การวิเคราะห์สมการอดิศัยเชิงเรขาคณิตและทฤษฎีบทการมีอยู่ของสถานะผูกพัน}
เพื่อแก้สมการอดิศัย (Transcendental Equations) ทั้งสอง เรานิยามตัวแปรไร้มิติ
\begin{equation}
\xi \equiv ka, \qquad \eta \equiv \kappa a, \qquad \xi, \eta > 0
\end{equation}
เมื่อนำกำลังสองของทั้งสองตัวแปรมารวมกัน จะได้สมการวงกลมที่มีรัศมี $R$ ซึ่งขึ้นอยู่กับมิติและความลึกของบ่อศักย์เท่านั้น
\begin{equation}
\xi^2 + \eta^2 = \frac{2m(E + V_0)a^2}{\hbar^2} + \frac{2m(-E)a^2}{\hbar^2} = \frac{2m V_0 a^2}{\hbar^2} \equiv R^2
\end{equation}

ดังนั้น การหาระดับพลังงานจึงเทียบเท่ากับการหาจุดตัดในควอดรันต์ที่หนึ่ง ($\xi > 0, \eta > 0$) ระหว่างเส้นรอบวงกลม $\xi^2 + \eta^2 = R^2$ กับเส้นโค้งอดิศัย
\begin{align}
\text{สถานะพาริตีคู่} \quad &\colon \quad \eta = \xi \tan\xi \label{eq_finite_even} \\
\text{สถานะพาริตีคี่} \quad &\colon \quad \eta = -\xi \cot\xi \label{eq_finite_odd}
\end{align}

\begin{figure}[H]
\centering
\begin{tikzpicture}[scale=1.2]
    % Coordinate axes
    \draw[-{Stealth}, darkNavy, line width=1.2pt] (0, 0) -- (5.5, 0) node[right] {$\xi = ka$};
    \draw[-{Stealth}, darkNavy, line width=1.2pt] (0, 0) -- (0, 4.5) node[above] {$\eta = \kappa a$};
    
    % Circle of radius R
    \draw[slateGrey, line width=1.5pt, dashed] (0, 3.8) arc (90:0:3.8);
    \node at (2.9, 2.9) [font=\small\bfseries\color{slateGrey}] {$R = \sqrt{2mV_0 a^2/\hbar^2}$};
    
    % Even curve eta = xi * tan(xi)
    \draw[cyberBlue, line width=1.8pt, domain=0:1.35, samples=60] plot (\x, {\x*tan(\x r)});
    \draw[cyberBlue, line width=1.8pt, domain=3.1415:4.2, samples=60] plot (\x, {\x*tan(\x r)});
    \node at (1.1, 3.8) [font=\small\bfseries\color{cyberBlue}] {คู่ $\eta = \xi\tan\xi$};
    
    % Odd curve eta = -xi * cot(xi)
    \draw[rubyRed, line width=1.8pt, domain=1.57:2.8, samples=60] plot (\x, {-\x/tan(\x r)});
    \node at (2.2, 4.2) [font=\small\bfseries\color{rubyRed}] {คี่ $\eta = -\xi\cot\xi$};
    
    % Asymptotes
    \draw[dotted, slateGrey] (1.57, 0) -- (1.57, 4.3) node[above, font=\footnotesize] {$\pi/2$};
    \draw[dotted, slateGrey] (3.14, 0) -- (3.14, 4.3) node[above, font=\footnotesize] {$\pi$};
    \draw[dotted, slateGrey] (4.71, 0) -- (4.71, 4.3) node[above, font=\footnotesize] {$3\pi/2$};
    
    % Intersection points
    \fill[goldAccent] (1.12, 2.3) circle (3pt) node[right, font=\footnotesize\bfseries] {สถานะคู่ที่ 1};
    \fill[goldAccent] (2.45, 2.85) circle (3pt) node[above right, font=\footnotesize\bfseries] {สถานะคี่ที่ 1};
\end{tikzpicture}
\caption{การหาผลเฉลยของสมการอดิศัยสำหรับบ่อศักย์จำกัดหนึ่งมิติด้วยวิธีเรขาคณิต}
\label{fig_transcendental_finite_well}
\end{figure}

จากกราฟในรูปที่ \ref{fig_transcendental_finite_well} นำไปสู่ข้อสรุปทางฟิสิกส์อันลึกซึ้งดังนี้
\begin{enumerate}
    \item \textbf{การมีอยู่เสมอของสถานะพื้น} ---- สำหรับสถานะพาริตีคู่ เมื่อ $\xi \to 0$ เส้นโค้ง $\xi\tan\xi \to 0$ ขณะที่เส้นรอบวงกลมเริ่มต้นที่ $\eta = R > 0$ ดังนั้น เส้นรอบวงกลมจะต้องตัดกับเส้นโค้ง $\xi\tan\xi$ อย่างน้อยหนึ่งจุดเสมอ ไม่ว่าค่า $R$ จะเล็กเพียงใดก็ตาม นั่นคือ \textbf{ในหนึ่งมิติ บ่อศักย์จำกัดจะมีสถานะผูกพันอย่างน้อยหนึ่งสถานะเสมอ} แม้บ่อจะตื้นหรือแคบเพียงใดก็ตาม
    
    \item \textbf{เงื่อนไขการเกิดสถานะคี่} ---- เส้นโค้งสถานะคี่ $\eta = -\xi\cot\xi$ จะเริ่มปรากฏค่าจริงที่แกน $\xi$ ณ จุด $\xi = \pi/2$ ดังนั้น สถานะพาริตีคี่สถานะแรกจะเกิดขึ้นได้ก็ต่อเมื่อรัศมีวงกลมมีค่ามากกว่า $\pi/2$ นั่นคือ
    \begin{equation}
    R > \frac{\pi}{2} \quad \Longrightarrow \quad V_0 a^2 > \frac{\pi^2 \hbar^2}{8m}
    \end{equation}
    
    \item \textbf{จำนวนสถานะผูกพันทั้งหมด} ---- จำนวนสถานะผูกพันทั้งหมด $N$ ในบ่อศักย์จำกัดหนึ่งมิติถูกกำหนดโดยรัศมี $R$ อย่างแม่นยำ
    \begin{equation}
    N = 1 + \left\lfloor \frac{2R}{\pi} \right\rfloor = 1 + \left\lfloor \frac{2a}{\pi\hbar}\sqrt{2m V_0} \right\rfloor
    \end{equation}
\end{enumerate}

% ==========================================================
% 4.5 บ่อศักย์และกำแพงศักย์เดลตาของดิแรก
% ==========================================================
\section{บ่อศักย์และกำแพงศักย์เดลตาของดิแรก}
\label{sec_delta_potential}
\index{ศักย์เดลตาของดิแรก (Dirac Delta Potential)}
\index{ความไม่ต่อเนื่องของอนุพันธ์ (Derivative Discontinuity)}
\index{แบบจำลองพันธะเคมี (Chemical Bonding Model)}

ฟังก์ชันเดลตาของดิแรก $\delta(x)$ เป็นการประมาณค่าในอุดมคติของสนามศักย์ที่มีอาณาบริเวณแคบมากแต่มีความเข้มข้นสูงยิ่งยวด เช่น รอยต่อระหว่างแผ่นตัวนำ หรือแรงนิวเคลียร์ระยะสั้น

\subsection{การอินทิเกรตข้ามจุดเอกฐานและความไม่ต่อเนื่องของอนุพันธ์}
พิจารณาสมการชเรอดิงเงอร์ที่มีศักย์เดลตา ณ ตำแหน่ง $x = 0$
\begin{equation}
-\frac{\hbar^2}{2m} \frac{d^2\psi}{dx^2} - \alpha \delta(x) \psi(x) = E\psi(x)
\end{equation}
โดยที่ $\alpha > 0$ แทนบ่อศักย์ดึงดูด และ $\alpha < 0$ แทนกำแพงศักย์ผลัก

เนื่องจากศักย์ลู่ออกสู่ค่าอนันต์ที่ $x = 0$ อนุพันธ์อันดับหนึ่งของฟังก์ชันคลื่นจึงไม่ต่อเนื่องที่จุดดังกล่าว ในการหาเงื่อนไขการต่อเชื่อม เราทำการอินทิเกรตสมการชเรอดิงเงอร์ในช่วงแคบๆ ข้ามจุดเอกฐานจาก $-\epsilon$ ถึง $+\epsilon$
\begin{equation}
-\frac{\hbar^2}{2m} \int_{-\epsilon}^{+\epsilon} \frac{d^2\psi}{dx^2} dx - \alpha \int_{-\epsilon}^{+\epsilon} \delta(x)\psi(x) dx = E \int_{-\epsilon}^{+\epsilon} \psi(x) dx
\end{equation}
เมื่อใช้คุณสมบัติการคัดกรองของฟังก์ชันเดลตา $\int_{-\epsilon}^{+\epsilon} \delta(x)\psi(x) dx = \psi(0)$ และเนื่องจากฟังก์ชันคลื่นต้องมีความต่อเนื่องที่จุดกำเนิด อินทิกรัลทางขวามือจะลู่เข้าสู่ศูนย์เมื่อ $\epsilon \to 0$
\begin{equation}
\lim_{\epsilon \to 0} \left[ \left.\frac{d\psi}{dx}\right|_{+\epsilon} - \left.\frac{d\psi}{dx}\right|_{-\epsilon} \right] = -\frac{2m\alpha}{\hbar^2}\psi(0)
\end{equation}
สมการนี้แสดงให้เห็นว่า \textbf{อนุพันธ์อันดับหนึ่งของฟังก์ชันคลื่นเกิดการกระโดดแบบไม่ต่อเนื่อง (Step Discontinuity)} เป็นขนาดที่แปรผันตรงกับค่าของฟังก์ชันคลื่น ณ จุดเอกฐาน

\subsection{สถานะผูกพันในบ่อศักย์เดลตา ($E < 0$)}
สำหรับสถานะผูกพัน $E = -\frac{\hbar^2\kappa^2}{2m} < 0$ ในบริเวณ $x \ne 0$ พลังงานศักย์เป็นศูนย์ ฟังก์ชันคลื่นที่ลู่เข้าสู่ศูนย์ที่ระยะอนันต์คือ
\begin{equation}
\psi(x) = \begin{cases}
A e^{\kappa x} & \text{เมื่อ } x < 0 \\
A e^{-\kappa x} & \text{เมื่อ } x > 0
\end{cases} \quad \Longrightarrow \quad \psi(x) = A e^{-\kappa |x|}
\end{equation}
คำนวณอนุพันธ์ทางซ้ายและขวาของจุดกำเนิด
\begin{equation}
\left.\frac{d\psi}{dx}\right|_{0^+} = -\kappa A, \qquad \left.\frac{d\psi}{dx}\right|_{0^-} = +\kappa A
\end{equation}
แทนค่าลงในเงื่อนไขความไม่ต่อเนื่องของอนุพันธ์
\begin{equation}
(-\kappa A) - (+\kappa A) = -\frac{2m\alpha}{\hbar^2} A \quad \Longrightarrow \quad -2\kappa = -\frac{2m\alpha}{\hbar^2} \quad \Longrightarrow \quad \kappa = \frac{m\alpha}{\hbar^2}
\end{equation}
ส่งผลให้บ่อศักย์เดลตาของดิแรกมี \textbf{สถานะผูกพันเพียงสถานะเดียวเท่านั้น} โดยมีระดับพลังงานและฟังก์ชันคลื่นที่นอร์มอลไลซ์แล้วดังนี้
\begin{align}
E &= -\frac{\hbar^2\kappa^2}{2m} = -\frac{m\alpha^2}{2\hbar^2} \\
\psi(x) &= \frac{\sqrt{m\alpha}}{\hbar} \exp\left(-\frac{m\alpha}{\hbar^2}|x|\right)
\end{align}

\subsection{แบบจำลองพันธะเคมีด้วยบ่อศักย์เดลตาคู่}
เมื่อขยายระบบไปสู่บ่อศักย์เดลตาสองตำแหน่ง $V(x) = -\alpha[\delta(x + a) + \delta(x - a)]$ เพื่อจำลองการแลกเปลี่ยนอิเล็กตรอนระหว่างสองนิวเคลียส (แบบจำลองอย่างง่ายของโมเลกุลไฮโดรเจนไอออน $H_2^+$) จะพบว่าสถานะผูกพันเกิดการแยกออกเป็นสองสถานะ
\begin{align}
\text{สถานะพันธะ (Bonding State, สมมาตรคู่)} \quad &\colon \quad E_+ = -\frac{\hbar^2\kappa_+^2}{2m}, \quad \kappa_+ = \frac{m\alpha}{\hbar^2}(1 + e^{-2\kappa_+ a}) \\
\text{สถานะปฏิพันธะ (Antibonding State, สมมาตรคี่)} \quad &\colon \quad E_- = -\frac{\hbar^2\kappa_-^2}{2m}, \quad \kappa_- = \frac{m\alpha}{\hbar^2}(1 - e^{-2\kappa_- a})
\end{align}
เนื่องจาก $\kappa_+ > \kappa_-$ ระดับพลังงานของสถานะพันธะ $E_+$ จึงอยู่ต่ำกว่าสถานะปฏิพันธะ $E_-$ เสมอ พลังงานที่ลดต่ำลงนี้อธิบายกลไกควอนตัมของการเกิด \textbf{พันธะโคเวเลนต์ (Covalent Bond)} ซึ่งเกิดจากการที่ฟังก์ชันคลื่นของอิเล็กตรอนสามารถแทรกซึมและแทรกสอดเสริมกันในช่องว่างระหว่างนิวเคลียสทั้งสอง

% ==========================================================
% 4.6 ปรากฏการณ์ควอนตัมทันเนลลิงผ่านกำแพงศักย์และการสั่นพ้อง
% ==========================================================
\section{ปรากฏการณ์ควอนตัมทันเนลลิงผ่านกำแพงศักย์และการสั่นพ้อง}
\label{sec_tunneling_resonances}
\index{ควอนตัมทันเนลลิง (Quantum Tunneling)}
\index{กำแพงศักย์สี่เหลี่ยม (Rectangular Barrier)}
\index{ปรากฏการณ์แรมเซาเออร์-ทาวน์เซนด์ (Ramsauer-Townsend Effect)}
\index{ไดโอดทันเนลลิงสั่นพ้อง (Resonant Tunneling Diode)}

ปรากฏการณ์ควอนตัมทันเนลลิง (Quantum Tunneling) คือกระบวนการที่อนุภาคสามารถทะลุผ่านกำแพงพลังงานศักย์ที่มีความสูงมากกว่าพลังงานรวมของอนุภาค ($E < V_0$) ซึ่งเป็นปรากฏการณ์ที่เป็นไปไม่ได้อย่างเด็ดขาดตามหลักกลศาสตร์คลาสสิก

\subsection{การอนุพัทธ์สัมประสิทธิ์การทะลุผ่านอย่างแม่นยำ}
พิจารณากำแพงศักย์สี่เหลี่ยมที่มีความสูง $V_0$ และความหนา $a$
\begin{equation}
V(x) = \begin{cases}
0 & \text{เมื่อ } x < 0 \quad (\text{บริเวณที่ I}) \\
V_0 & \text{เมื่อ } 0 \le x \le a \quad (\text{บริเวณที่ II}) \\
0 & \text{เมื่อ } x > a \quad (\text{บริเวณที่ III})
\end{cases}
\end{equation}

สำหรับอนุภาคที่มีพลังงาน $E < V_0$ ฟังก์ชันคลื่นในแต่ละบริเวณเขียนได้ดังนี้
\begin{align}
\psi_I(x) &= e^{ik x} + r e^{-ik x} \quad (x < 0) \\
\psi_{II}(x) &= A e^{-\kappa x} + B e^{\kappa x} \quad (0 \le x \le a) \\
\psi_{III}(x) &= t e^{ik x} \quad (x > a)
\end{align}
โดยที่ $k = \frac{\sqrt{2mE}}{\hbar}$ และ $\kappa = \frac{\sqrt{2m(V_0 - E)}}{\hbar}$

\begin{figure}[H]
\centering
\begin{tikzpicture}[scale=1.15]
    % Axes
    \draw[-{Stealth}, darkNavy, line width=1.2pt] (-0.8, 0) -- (6.8, 0) node[right] {$x$};
    \draw[-{Stealth}, darkNavy, line width=1.2pt] (0, -0.3) -- (0, 3.8) node[above] {$V(x), \psi(x)$};
    
    % Barrier
    \fill[slateGrey!25!white] (1.8, 0) rectangle (3.8, 2.8);
    \draw[darkNavy, line width=2.2pt] (-0.5, 0) -- (1.8, 0) -- (1.8, 2.8) -- (3.8, 2.8) -- (3.8, 0) -- (6.5, 0);
    \node at (2.8, 3.1) [font=\bfseries\color{darkNavy}] {กำแพงศักย์ $V_0$};
    
    % Energy line
    \draw[dashed, rubyRed, line width=1.3pt] (-0.5, 1.8) -- (6.5, 1.8) node[right, font=\bfseries\color{rubyRed}] {$E < V_0$};
    
    % Wavefunction Incident & Reflected (Region I)
    \draw[cyberBlue, line width=1.6pt, domain=-0.5:1.8, samples=80] plot (\x, {1.8 + 0.55*sin(12*(\x-1.8) r)});
    \node at (0.6, 2.6) [font=\small\bfseries\color{cyberBlue}] {บริเวณ I (ตกกระทบ/สะท้อน)};
    
    % Wavefunction Tunneling (Region II)
    \draw[quantumGold, line width=2pt, domain=1.8:3.8, samples=50] plot (\x, {1.8 + 0.55*exp(-1.3*(\x-1.8)) - 0.05*exp(1.3*(\x-3.8))});
    \node at (2.8, 1.2) [font=\footnotesize\bfseries\color{quantumGold}] {$\psi \sim e^{-\kappa x}$};
    
    % Wavefunction Transmitted (Region III)
    \draw[emeraldGlow, line width=1.6pt, domain=3.8:6.5, samples=80] plot (\x, {1.8 + 0.16*sin(12*(\x-3.8) r)});
    \node at (5.2, 2.6) [font=\small\bfseries\color{emeraldGlow}] {บริเวณ III (ทะลุผ่าน)};
\end{tikzpicture}
\caption{พฤติกรรมการลดทอนและทะลุผ่านของฟังก์ชันคลื่นในปรากฏการณ์ควอนตัมทันเนลลิง (เป้าหมายติดตามภาพ AR Image Target)}
\label{fig_tunneling_barrier_wave}
\end{figure}

การจับคู่เงื่อนไขความต่อเนื่องของ $\psi(x)$ และ $\frac{d\psi}{dx}$ ที่ตำแหน่ง $x = 0$ และ $x = a$ นำไปสู่ระบบสมการ 4 สมการ
\begin{align}
1 + r &= A + B \\
ik(1 - r) &= -\kappa A + \kappa B \\
A e^{-\kappa a} + B e^{\kappa a} &= t e^{ik a} \\
-\kappa A e^{-\kappa a} + \kappa B e^{\kappa a} &= ik t e^{ik a}
\end{align}

เมื่อแก้ระบบสมการเชิงพีชคณิตเพื่อหาแอมพลิจูดการทะลุผ่าน $t$ และคำนวณสัมประสิทธิ์การทะลุผ่าน $T = |t|^2$ จะได้ผลลัพธ์ที่แม่นยำสมบูรณ์คือ
\begin{equation}
T = \frac{1}{1 + \frac{V_0^2}{4E(V_0 - E)} \sinh^2(\kappa a)} \label{eq_tunneling_exact}
\end{equation}

ในขีดจำกัดที่กำแพงศักย์มีความหนามากหรืออนุภาคมีพลังงานต่ำกว่ายอดกำแพงอย่างชัดเจน ($\kappa a \gg 1$) ฟังก์ชันไฮเพอร์โบลิกไซน์สามารถประมาณค่าได้ด้วย $\sinh(\kappa a) \approx \frac{1}{2}e^{\kappa a}$ ส่งผลให้สมการ \eqref{eq_tunneling_exact} ลดรูปลงเป็น
\begin{equation}
T \approx 16 \frac{E}{V_0}\left(1 - \frac{E}{V_0}\right) e^{-2\kappa a} \approx e^{-2\kappa a}
\end{equation}
สมการนี้แสดงให้เห็นว่า \textbf{ความน่าจะเป็นในการทะลุกำแพงศักย์ลดลงอย่างรวดเร็วแบบเอกซ์โพเนนเชียลตามความหนาของกำแพง $a$ และสแควร์รูทของความสูงกำแพง $\sqrt{V_0 - E}$}

\subsection{การทะลุผ่านเมื่อ $E > V_0$ และปรากฏการณ์แรมเซาเออร์-ทาวน์เซนด์}
เมื่ออนุภาคมีพลังงานสูงกว่ากำแพงศักย์ ($E > V_0$) เราสามารถแปลงผลเฉลยข้างต้นได้โดยตรงโดยแทนที่ $\kappa \to i k_2$ โดยที่ $k_2 = \frac{\sqrt{2m(E - V_0)}}{\hbar}$ และใช้เอกลักษณ์ $\sinh(i k_2 a) = i \sin(k_2 a)$ จะได้
\begin{equation}
T = \frac{1}{1 + \frac{V_0^2}{4E(E - V_0)} \sin^2(k_2 a)}
\end{equation}

ข้อค้นพบที่น่าทึ่งคือ เมื่อความหนาของกำแพงสอดคล้องกับเงื่อนไข
\begin{equation}
k_2 a = n\pi \quad \Longrightarrow \quad a = \frac{n\lambda_2}{2}, \qquad n = 1, 2, 3, \dots
\end{equation}
เทอม $\sin^2(k_2 a)$ จะกลายเป็นศูนย์ ส่งผลให้สัมประสิทธิ์การทะลุผ่านมีค่าเท่ากับหนึ่งอย่างสมบูรณ์ ($T = 1$) อนุภาคจะเคลื่อนที่ทะลุผ่านกำแพงศักย์ไปได้โดยไม่มีการสะท้อนกลับแม้แต่น้อย ปรากฏการณ์นี้เรียกว่า \textbf{การสั่นพ้องของการทะลุผ่าน (Transmission Resonance)} ซึ่งอธิบายการทดลองของคาร์ล แรมเซาเออร์ (Carl Ramsauer) และจอห์น ทาวน์เซนด์ (John Townsend) ในปี ค.ศ. 1921 ที่พบว่าก๊าซเฉื่อย (เช่น ซีนอนและคริปทอน) มีความโปร่งใสต่อลำอิเล็กตรอนที่มีพลังงานจลน์จำเพาะบางค่าอย่างสมบูรณ์แบบ

\subsection{ไดโอดทันเนลลิงสั่นพ้อง (Resonant Tunneling Diode หรือ RTD)}
เมื่อนำกำแพงศักย์สองชั้นมาวางขนานกันโดยเว้นระยะห่างตรงกลาง จะเกิดโครงสร้างบ่อควอนตัมคั่นกลางระหว่างกำแพงศักย์คู่ (Double Barrier Quantum Well)

\begin{figure}[H]
\centering
\begin{tikzpicture}[scale=1.1]
    % Profile
    \draw[-{Stealth}, darkNavy, line width=1.2pt] (-0.5, 0) -- (7.5, 0) node[right] {$x$};
    \draw[-{Stealth}, darkNavy, line width=1.2pt] (0, -0.2) -- (0, 3.5) node[above] {$V(x)$};
    
    % Barrier 1 and 2
    \fill[slateGrey!30] (1.2, 0) rectangle (2.2, 2.8);
    \fill[slateGrey!30] (4.0, 0) rectangle (5.0, 2.8);
    \draw[darkNavy, line width=2pt] (0, 0) -- (1.2, 0) -- (1.2, 2.8) -- (2.2, 2.8) -- (2.2, 0) -- (4.0, 0) -- (4.0, 2.8) -- (5.0, 2.8) -- (5.0, 0) -- (7.2, 0);
    
    % Resonant state in well
    \draw[dashed, rubyRed, line width=1.5pt] (2.2, 1.3) -- (4.0, 1.3) node[midway, above, font=\bfseries\color{rubyRed}] {ระดับสั่นพ้อง $E_R$};
    
    % Tunneling arrows
    \draw[-{Stealth}, cyberBlue, line width=1.8pt] (0.2, 1.3) -- (1.1, 1.3);
    \draw[-{Stealth}, emeraldGlow, line width=1.8pt] (5.1, 1.3) -- (6.5, 1.3) node[right, font=\bfseries\color{emeraldGlow}] {$T \approx 1$};
\end{tikzpicture}
\caption{โครงสร้างกำแพงศักย์คู่และระดับพลังงานสั่นพ้องในไดโอด RTD}
\label{fig_rtd_profile}
\end{figure}

หากพลังงานของอิเล็กตรอนตกกระทบตรงกับระดับพลังงานกึ่งผูกพัน (Quasi-bound State, $E_R$) ภายในบ่อศักย์ คลื่นจะเกิดการแทรกสอดเสริมกันอย่างรุนแรงภายในบ่อ ส่งผลให้สัมประสิทธิ์การทะลุผ่านพุ่งสูงขึ้นสู่ค่าเกือบ $100\%$ ตามสูตรไบรต์-วิกเนอร์ (Breit-Wigner Formula)
\begin{equation}
T(E) \approx \frac{(\Gamma/2)^2}{(E - E_R)^2 + (\Gamma/2)^2}
\end{equation}
โดยที่ $\Gamma$ คือความกว้างของเรโซแนนซ์ที่สัมพันธ์กับอายุขัยของสถานะกึ่งผูกพัน $\tau = \hbar/\Gamma$

เมื่อป้อนแรงดันไฟฟ้าไบแอส $V_{\text{bias}}$ ให้แก่อุปกรณ์ ระดับพลังงาน $E_R$ จะถูกเลื่อนลงมาให้ตรงกับแถบการนำไฟฟ้าของขั้วส่ง ทำให้กระแสไฟฟ้าพุ่งขึ้นสู่ค่าสูงสุด (Peak Current) แต่เมื่อเพิ่มแรงดันไฟฟ้าต่อไป ระดับ $E_R$ จะเลื่อนหลุดพ้นจากแถบการนำไฟฟ้า ส่งผลให้กระแสไฟฟ้าลดวูบลงทันที ก่อให้เกิดสมบัติ \textbf{ความต้านทานเชิงอนุพันธ์เป็นลบ (Negative Differential Resistance หรือ NDR)} ซึ่งนำไปใช้สร้างวงจรกำเนิดสัญญาณความถี่สูงยิ่งยวดในย่านเทระเฮิร์ตซ์ (Terahertz Oscillators) ที่ความถี่สูงกว่า 1.5 THz ในอุปกรณ์สื่อสาร 6G

% ==========================================================
% 4.7 การประมาณค่า WKB และทฤษฎีแอลฟาของกามอฟ
% ==========================================================
\section{การประมาณค่า WKB และทฤษฎีแอลฟาของกามอฟ}
\label{sec_wkb_gamow}
\index{การประมาณค่า WKB (WKB Approximation)}
\index{จุดวกกลับคลาสสิก (Classical Turning Point)}
\index{การสลายตัวให้อนุภาคแอลฟา (Alpha Decay)}
\index{กฎไกเกอร์-นัททอลล์ (Geiger-Nuttall Law)}

เมื่อฟังก์ชันพลังงานศักย์ $V(x)$ มีรูปร่างที่แปรเปลี่ยนไปอย่างต่อเนื่องตามตำแหน่ง สมการชเรอดิงเงอร์จะไม่สามารถหาผลเฉลยแม่นยำในรูปฟังก์ชันพื้นฐานได้ การประมาณค่าแบบเซมิคลาสสิกที่พัฒนาโดยเกรเกอร์ เวนต์เซล (Gregor Wentzel), เฮนดริก คราเมอร์ส (Hendrik Kramers) และเลออน บริลลูแอง (Léon Brillouin) ในปี ค.ศ. 1926 หรือ \textbf{การประมาณค่า WKB} จึงเป็นเครื่องมือทรงพลังในการคำนวณสัมประสิทธิ์การทะลุผ่าน

\subsection{สมมติฐานเซมิคลาสสิกและฟังก์ชันคลื่น WKB}
ภายใต้เงื่อนไขที่ความยาวคลื่นเดอบรอยล์ของอนุภาคเปลี่ยนแปลงไปอย่างช้าๆ เทียบกับระยะทาง ($\left|\frac{d\lambda}{dx}\right| \ll 1$) เราตั้งสมมติฐานฟังก์ชันคลื่นในรูป
\begin{equation}
\psi(x) = \exp\left[ \frac{i}{\hbar} S(x) \right], \qquad S(x) = S_0(x) + \hbar S_1(x) + \dots
\end{equation}
เมื่อกระจายตามเลขชี้กำลังของ $\hbar$ และจัดรูป จะได้ฟังก์ชันคลื่น WKB ในสองบริเวณดังนี้
\begin{align}
\text{บริเวณที่อนุญาตคลาสสิก } (E > V(x)) \quad &\colon \quad \psi(x) \approx \frac{C}{\sqrt{p(x)}} \exp\left[ \pm \frac{i}{\hbar}\int p(x') dx' \right] \\
\text{บริเวณต้องห้ามคลาสสิก } (E < V(x)) \quad &\colon \quad \psi(x) \approx \frac{D}{\sqrt{|p(x)|}} \exp\left[ \pm \frac{1}{\hbar}\int |p(x')| dx' \right]
\end{align}
โดยที่โมเมนตัมคลาสสิกกำหนดโดย $p(x) = \sqrt{2m(E - V(x))}$

สัมประสิทธิ์การทะลุผ่านของกำแพงศักย์รูปทรงใดๆ ที่มีจุดวกกลับคลาสสิก (Classical Turning Points) ณ ตำแหน่ง $x_1$ และ $x_2$ คำนวณได้จาก \textbf{อินทิกรัลกามอฟ (Gamow Factor)}
\begin{equation}
T \approx \exp\left[ -2 \gamma \right] = \exp\left[ -\frac{2}{\hbar} \int_{x_1}^{x_2} \sqrt{2m(V(x) - E)} \, dx \right]
\end{equation}

\subsection{ทฤษฎีการสลายตัวให้อนุภาคแอลฟาของจอร์จ กามอฟ}
ในคริสต์ทศวรรษ 1920 ปริศนาที่ยิ่งใหญ่ที่สุดของฟิสิกส์นิวเคลียร์คือ อนุภาคแอลฟา ($\alpha$, นิวเคลียสของฮีเลียม-4) ที่ปล่อยออกมาจากการสลายตัวของยูเรเนียม-238 มีพลังงานจลน์เพียง $E_\alpha \approx 4.2\text{ MeV}$ ทว่าการทดลองยิงอนุภาคแอลฟากลับเข้าไปในนิวเคลียสของเออร์เนสต์ รัทเทอร์ฟอร์ด แสดงให้เห็นว่ากำแพงศักย์คูลอมบ์ของนิวเคลียสสูงถึง $V_C \approx 25\text{ MeV}$ ตามฟิสิกส์คลาสสิก อนุภาคแอลฟาไม่มีวันที่จะหลุดรอดออกมาได้

ในปี ค.ศ. 1928 จอร์จ กามอฟ (George Gamow) รวมถึงโรนัลด์ เกอร์นีย์ (Ronald Gurney) และเอ็ดเวิร์ด คอนดอน (Edward Condon) ได้เสนอว่าการสลายตัวของนิวเคลียสเกิดจาก \textbf{ควอนตัมทันเนลลิง} ของอนุภาคแอลฟาผ่านกำแพงศักย์คูลอมบ์

\begin{figure}[H]
\centering
\begin{tikzpicture}[scale=1.1]
    % Axes
    \draw[-{Stealth}, darkNavy, line width=1.2pt] (0, 0) -- (6.8, 0) node[right] {$r$};
    \draw[-{Stealth}, darkNavy, line width=1.2pt] (0, -1.8) -- (0, 3.8) node[above] {$V(r)$};
    
    % Nuclear Well
    \draw[darkNavy, line width=2pt] (0, -1.5) -- (1.5, -1.5) -- (1.5, 3.2);
    \node at (0.75, -1.0) [font=\small\bfseries\color{darkNavy}] {บ่อนิวเคลียร์};
    
    % Coulomb Barrier
    \draw[darkNavy, line width=2pt, domain=1.5:6.2, samples=60] plot (\x, {4.8/\x});
    \node at (3.5, 2.2) [font=\small\bfseries\color{darkNavy}] {$V(r) = \frac{2Ze^2}{4\pi\varepsilon_0 r}$};
    
    % Energy line
    \draw[dashed, rubyRed, line width=1.4pt] (0, 1.2) -- (6.2, 1.2) node[right, font=\bfseries\color{rubyRed}] {$E_\alpha$};
    
    % Turning points
    \draw[dotted, slateGrey] (1.5, 0) -- (1.5, 1.2) node[below=14pt, font=\footnotesize\bfseries] {$R_0$};
    \draw[dotted, slateGrey] (4.0, 0) -- (4.0, 1.2) node[below=14pt, font=\footnotesize\bfseries] {$r_c$};
    
    % Shaded Tunneling Region
    \fill[slateGrey!30, domain=1.5:4.0, variable=\x] (1.5, 1.2) -- plot (\x, {4.8/\x}) -- (4.0, 1.2) -- cycle;
    \node at (2.7, 1.5) [font=\footnotesize\bfseries\color{quantumGold}] {ทันเนลลิง};
\end{tikzpicture}
\caption{กำแพงศักย์คูลอมบ์และกลไกควอนตัมทันเนลลิงในการสลายตัวให้อนุภาคแอลฟาของกามอฟ}
\label{fig_gamow_alpha_decay}
\end{figure}

ศักย์ไฟฟ้าคูลอมบ์ภายนอกรัศมีนิวเคลียส $R_0$ กำหนดโดย
\begin{equation}
V(r) = \frac{1}{4\pi\varepsilon_0} \frac{2(Z - 2)e^2}{r} \equiv \frac{2 Z' e^2}{4\pi\varepsilon_0 r}
\end{equation}
จุดวกกลับคลาสสิกภายนอกคือ $r_c$ ซึ่ง $V(r_c) = E_\alpha \implies r_c = \frac{2 Z' e^2}{4\pi\varepsilon_0 E_\alpha}$

คำนวณเลขชี้กำลังของกามอฟ $2\gamma$
\begin{equation}
2\gamma = \frac{2}{\hbar} \int_{R_0}^{r_c} \sqrt{2m_\alpha \left( \frac{2Z' e^2}{4\pi\varepsilon_0 r} - E_\alpha \right)} dr = \frac{2\sqrt{2m_\alpha E_\alpha}}{\hbar} \int_{R_0}^{r_c} \sqrt{\frac{r_c}{r} - 1} \, dr
\end{equation}
เปลี่ยนตัวแปร $r = r_c \cos^2\theta$ จะได้ผลการอินทิเกรต
\begin{equation}
2\gamma = \frac{4 Z' e^2}{4\pi\varepsilon_0 \hbar v_\alpha} \left[ \arccos\left(\sqrt{\frac{R_0}{r_c}}\right) - \sqrt{\frac{R_0}{r_c}\left(1 - \frac{R_0}{r_c}\right)} \right]
\end{equation}
เนื่องจาก $r_c \gg R_0$ สามารถประมาณค่าในวงเล็บได้ด้วย $\frac{\pi}{2} - 2\sqrt{\frac{R_0}{r_c}}$ นำไปสู่
\begin{equation}
2\gamma \approx \frac{2\pi Z' e^2}{4\pi\varepsilon_0 \hbar v_\alpha} - \frac{8}{\hbar}\sqrt{\frac{Z' e^2 m_\alpha R_0}{4\pi\varepsilon_0}} \equiv C_1 \frac{Z'}{\sqrt{E_\alpha}} - C_2 \sqrt{Z' R_0}
\end{equation}

อัตราการสลายตัว $\lambda$ สัมพันธ์กับความถี่ที่อนุภาคแอลฟาวิ่งชนผนังนิวเคลียส $\nu_0 = \frac{v_{\text{in}}}{2R_0} \approx 10^{21}\text{ s}^{-1}$
\begin{equation}
\lambda = \nu_0 T = \nu_0 e^{-2\gamma} \quad \Longrightarrow \quad \ln\lambda = \ln\nu_0 - C_1 \frac{Z'}{\sqrt{E_\alpha}} + C_2 \sqrt{Z' R_0}
\end{equation}
เมื่อแปลงเป็นครึ่งชีวิต $t_{1/2} = \frac{\ln 2}{\lambda}$ จะได้ \textbf{กฎไกเกอร์-นัททอลล์ (Geiger-Nuttall Law)}
\begin{equation}
\log_{10} t_{1/2} = A \frac{Z}{\sqrt{E_\alpha}} + B
\end{equation}
ทฤษฎีนี้ประสบความสำเร็จอย่างงดงามในการอธิบายความจริงเชิงประจักษ์ที่ว่า \textbf{เหตุใดพลังงานของอนุภาคแอลฟาที่เพิ่มขึ้นเพียง 2 เท่า (จาก 4 MeV เป็น 8.8 MeV) จึงส่งผลให้ครึ่งชีวิตของนิวเคลียสลดลงอย่างมหาศาลถึง $10^{24}$ เท่า} (จาก $4.5 \times 10^9$ ปีใน $^{238}\text{U}$ เหลือเพียง $0.3\times 10^{-6}$ วินาทีใน $^{212}\text{Po}$)

% ==========================================================
% 4.8 การประยุกต์ล้ำยุค: กล้อง STM สเปกโทรสโกปี STS และชิปซับนาโนเมตร
% ==========================================================
\section{การประยุกต์ล้ำยุค กล้อง STM สเปกโทรสโกปี STS และชิปซับนาโนเมตร}
\label{sec_modern_stm_nano}
\index{กล้องจุลทรรศน์ส่องกราดแบบอุโมงค์ (STM)}
\index{ทฤษฎีเทอร์ซอฟฟ์-ฮามันน์ (Tersoff-Hamann Theory)}
\index{สเปกโทรสโกปีแบบอุโมงค์ (STS)}
\index{ทรานซิสเตอร์ RibbonFET (RibbonFET GAA)}

ความเข้าใจอย่างลึกซึ้งในกลศาสตร์ควอนตัมทันเนลลิงได้เปลี่ยนผ่านจากการทดลองทางทฤษฎีไปสู่อภิมหาเทคโนโลยีระดับนาโนเมตรที่ขับเคลื่อนอารยธรรมมนุษย์ในปัจจุบัน

\subsection{ทฤษฎีเทอร์ซอฟฟ์-ฮามันน์สำหรับกล้องจุลทรรศน์ STM}
เกิร์ด บินนิก (Gerd Binnig) และไฮน์ริช โรเรอร์ (Heinrich Rohrer) ได้รับรางวัลโนเบลสาขาฟิสิกส์ในปี ค.ศ. 1986 จากการประดิษฐ์ \textbf{กล้องจุลทรรศน์ส่องกราดแบบอุโมงค์ (Scanning Tunneling Microscope หรือ STM)}

\begin{figure}[H]
\centering
\begin{tikzpicture}[scale=1.1]
    % Tip
    \fill[darkNavy!70!white] (2.8, 3.8) -- (3.2, 3.8) -- (3.05, 2.5) -- (2.95, 2.5) -- cycle;
    \fill[quantumGold] (3.0, 2.4) circle (4pt);
    \node at (3.0, 2.4) [right=6pt, font=\footnotesize\bfseries\color{quantumGold}] {อะตอมปลายเข็ม};
    \node at (3.0, 4.1) [font=\small\bfseries\color{darkNavy}] {หัวเข็มทังสเตน/แพลทินัม};
    
    % Surface atoms
    \foreach \x in {0.5, 1.3, 2.1, 2.9, 3.7, 4.5, 5.3} {
        \fill[cyberBlue!70!white] (\x, 1.0) circle (8pt);
        \fill[cyberBlue] (\x, 0.4) circle (8pt);
    }
    \node at (6.3, 0.8) [font=\small\bfseries\color{cyberBlue}] {พื้นผิวตัวอย่าง};
    
    % Tunneling gap and current
    \draw[{Stealth}-{Stealth}, rubyRed, line width=1.5pt] (2.5, 1.4) -- (2.5, 2.3) node[midway, left, font=\footnotesize\bfseries] {$d \approx 0.5\text{ nm}$};
    \draw[-{Stealth}, dashed, rubyRed, line width=1.5pt] (3.0, 2.2) -- (3.0, 1.4) node[midway, right, font=\footnotesize\bfseries] {$I_T \propto e^{-2\kappa d}$};
    
    % Bias Voltage
    \draw[line width=1.2pt] (3.0, 3.8) -- (3.0, 4.5) -- (6.5, 4.5) -- (6.5, 1.0) -- (5.7, 1.0);
    \node at (6.5, 2.8) [draw, fill=white, font=\footnotesize\bfseries] {แรงดันไบแอส $V_B$};
\end{tikzpicture}
\caption{แผนภาพการทำงานของกล้องจุลทรรศน์ STM และการไหลของกระแสทันเนลลิงระดับอะตอมเดี่ยว}
\label{fig_stm_atomic_schematic}
\end{figure}

ตามทฤษฎีของเจอร์รี เทอร์ซอฟฟ์ (Jerry Tersoff) และดี. อาร์. ฮามันน์ (D. R. Hamann) ที่ตีพิมพ์ในปี ค.ศ. 1983 และ 1985 กระแสทันเนลลิง $I$ ภายใต้แรงดันไบแอสต่ำ $V$ และอุณหภูมิต่ำคำนวณได้จากทฤษฎีการรบกวนของบาร์ดีน (Bardeen's Transfer Hamiltonian Formalism)
\begin{equation}
I(V, \mathbf{r}_0) \propto e V \sum_{\nu} |\psi_\nu(\mathbf{r}_0)|^2 \delta(E_\nu - E_F) = e V \rho_s(\mathbf{r}_0, E_F)
\end{equation}
โดยที่ $\rho_s(\mathbf{r}_0, E_F)$ คือ \textbf{ความหนาแน่นสถานะเฉพาะที่ (Local Density of States หรือ LDOS)} ของพื้นผิวตัวอย่าง ณ ระดับพลังงานแฟร์มี $E_F$ วัดที่ตำแหน่งจุดศูนย์กลางความโค้งของปลายเข็ม $\mathbf{r}_0$

เนื่องจากฟังก์ชันคลื่นสลายตัวแบบเอกซ์โพเนนเชียลในช่องว่างสุญญากาศ กระแสทันเนลลิงจะขึ้นอยู่กับระยะห่าง $d$ อย่างรุนแรง
\begin{equation}
I \propto \exp\left( -2\kappa d \right) = \exp\left( -\frac{2\sqrt{2m\Phi}}{\hbar} d \right)
\end{equation}
โดยที่ $\Phi$ คือฟังก์ชันงานเฉลี่ย (Work Function, $\sim 4\text{ eV}$) เมื่อคำนวณตัวเลขจะพบว่า $\kappa \approx 10\text{ nm}^{-1}$ ส่งผลให้ \textbf{การเลื่อนปลายเข็มห่างออกไปเพียง $0.1\text{ nm}$ (ขนาดเส้นผ่านศูนย์กลางของอะตอมหนึ่งตัว) จะทำให้กระแสทันเนลลิงลดลงถึง $10$ เท่า} ความไวขั้นอัจฉริยะนี้ทำให้อุปกรณ์สามารถควบคุมความสูงได้ละเอียดถึง $0.001\text{ nm}$ และสามารถแยกแยะอะตอมเดี่ยวได้อย่างคมชัด

\subsection{สเปกโทรสโกปีแบบอุโมงค์ส่องกราด (Scanning Tunneling Spectroscopy หรือ STS)}
นอกจากการถ่ายภาพสามมิติของพื้นผิวแล้ว หากคงตำแหน่งปลายเข็มไว้กับที่แล้วปรับเปลี่ยนแรงดันไบแอส $V$ พร้อมวัดอนุพันธ์ของกระแสเทียบกับแรงดัน $\frac{dI}{dV}$ จะได้ข้อมูลสเปกโทรสโกปีทางอิเล็กทรอนิกส์
\begin{equation}
\frac{dI}{dV}(V) \propto \rho_s(\mathbf{r}_0, E_F + e V)
\end{equation}
เทคนิค STS เปิดโอกาสให้นักฟิสิกส์สามารถ \textbf{มองเห็นแถบพลังงาน ช่องว่างสภาพนำยิ่งยวด (Superconducting Gap) กรวยดิแรกในแกรฟีน และสถานะขอบโทโพโลยี (Majorana Bound States)} ได้โดยตรงในระดับอะตอมเดี่ยว ซึ่งเป็นรากฐานของการวิจัยควอนตัมคอมพิวเตอร์เชิงโทโพโลยีในปัจจุบัน

\subsection{วิศวกรรมชิปประมวลผลยุคซับนาโนเมตร จาก FinFET สู่ RibbonFET GAA}
ในอุตสาหกรรมเซมิคอนดักเตอร์ การย่อขนาดทรานซิสเตอร์ลงสู่ระดับต่ำกว่า 3 นาโนเมตรทำให้ความหนาของชั้นฉนวนเกตไดอิเล็กทริก ($SiO_2$ หรือ $HfO_2$) บางลงเหลือไม่ถึง $1\text{ nm}$ ส่งผลให้อิเล็กตรอนเกิด \textbf{กระแสไฟฟ้ารั่วไหลแบบทันเนลลิงโดยตรง (Direct Quantum Leakage Current)} อย่างรุนแรง ทำให้ทรานซิสเตอร์ไม่สามารถปิดสวิตช์ได้อย่างสมบูรณ์

เพื่อแก้ไขปัญหานี้ วิศวกรนาโนอิเล็กทรอนิกส์จึงได้พัฒนาสถาปัตยกรรมทรานซิสเตอร์ยุคใหม่ ได้แก่
\begin{enumerate}
    \item \textbf{High-$k$ Metal Gate (HKMG)} ---- การใช้วัสดุที่มีค่าคงที่ไดอิเล็กทริกสูง เช่น ฮาฟเนียมออกไซด์ ($HfO_2$) ทำให้สามารถเพิ่มความหนาเชิงกายภาพของชั้นฉนวนเพื่อลดการทันเนลลิงได้ โดยยังคงรักษาค่าความจุไฟฟ้าเกตไว้เท่าเดิม
    \item \textbf{Gate-All-Around (GAA) หรือ RibbonFET} ---- การสร้างแผ่นช่องนำกระแสขนาดนาโน (Nanosheets) ซ้อนกันในแนวดิ่งและโอบล้อมด้วยขั้วเกตทั้งสี่ด้าน ช่วยเพิ่มการควบคุมสนามไฟฟ้าเชิงคลาสสิกเพื่อต่อสู้กับผลกระทบของควอนตัมทันเนลลิงในชิปสถาปัตยกรรม 2 นาโนเมตร (Intel 18A และ TSMC N2)
\end{enumerate}

% ==========================================================
% 4.9 กล่องปฏิบัติการ WebAR & Image Tracking ประจำบทที่ 4
% ==========================================================
\section{ห้องปฏิบัติการเสมือนจริง WebAR ปรากฏการณ์ควอนตัมทันเนลลิงและกล้อง STM}
\label{sec_webar_lab04}

\begin{qrcodelabbox}{การทดลองควอนตัมทันเนลลิงและกล้อง STM ระดับอะตอม WebAR}{https://tsanaphysics.github.io/ar-Quantum/lab04_quantum_tunneling.html}
\begin{minipage}[c]{0.70\textwidth}
\textbf{คำแนะนำการทดลองเสมือนจริงสามมิติ (Interactive Tunneling Lab)}
\begin{enumerate}
    \item สแกนรหัสตอบสนองเร็ว (QR Code) ทางด้านขวา เพื่อเข้าสู่ห้องปฏิบัติการเสมือนจริง WebAR บนเว็บเบราว์เซอร์
    \item หันกล้องสมาร์ตโฟนเข้าหา \textbf{ภาพที่ \ref{fig_tunneling_barrier_wave}} เพื่อเรนเดอร์กำแพงศักย์ ฟังก์ชันคลื่น และหัวเข็ม STM สามมิติแบบลอยตัวเสมือนจริง
    \item ปรับเลื่อนความหนาของกำแพงศักย์ $a$ และความสูง $V_0$ เพื่อสังเกตการหดหายของฟังก์ชันคลื่นและการแปรผันของสัมประสิทธิ์การทะลุผ่าน $T$ แบบเรียลไทม์
    \item ทดลองขยับปลายเข็ม STM เหนือผิวหน้าผลึกทองคำจำลอง เพื่อดูการตอบสนองของกระแสไฟฟ้าทันเนลลิงตามระยะทาง $I \propto e^{-2\kappa d}$
\end{enumerate}
\end{minipage}
\hfill
\begin{minipage}[c]{0.25\textwidth}
\centering
\qrcode[height=2.8cm]{https://tsanaphysics.github.io/ar-Quantum/lab04_quantum_tunneling.html}
\vspace{4pt}
\scriptsize\textbf{สแกนเพื่อเปิด WebAR Lab 04}
\end{minipage}
\end{qrcodelabbox}

% ==========================================================
% 4.10 การจำลองเชิงตัวเลขด้วย Python: ระเบียบวิธี Transfer Matrix Method
% ==========================================================
\section{การจำลองเชิงตัวเลขด้วยภาษา Python ด้วยระเบียบวิธีเมทริกซ์ส่งผ่าน}
\label{sec_python_tmm}
\index{ระเบียบวิธีเมทริกซ์ส่งผ่าน (Transfer Matrix Method)}
\index{การจำลองด้วยภาษา Python (Python Quantum Simulation)}

ระเบียบวิธีเมทริกซ์ส่งผ่าน (Transfer Matrix Method หรือ TMM) เป็นขั้นตอนวิธีเชิงตัวเลขที่มีประสิทธิภาพสูงสุดในการคำนวณสัมประสิทธิ์การทะลุผ่านและการสะท้อนของอนุภาคควอนตัมผ่านโครงสร้างศักย์หนึ่งมิติที่มีรูปร่างซับซ้อนใดๆ โปรแกรมด้านล่างนี้จำลองการทะลุผ่านกำแพงศักย์คู่ในไดโอด RTD พร้อมแสดงการสั่นพ้องของระดับพลังงาน

\begin{lstlisting}[language=Python, caption={การจำลองการทะลุกำแพงศักย์คู่และไดโอด RTD ด้วย Transfer Matrix Method}, label={code_tmm_solver}]
import numpy as np
import matplotlib.pyplot as plt

# ค่าคงที่ทางฟิสิกส์ (หน่วยมาตรฐาน SI และอิเล็กตรอนโวลต์)
hbar = 1.054571817e-34       # จูล-วินาที
m_e = 9.1093837e-31          # มวลอิเล็กตรอน (กก.)
eV = 1.602176634e-19         # 1 eV ในหน่วยจูล

def transfer_matrix_barrier(E_val, x_grid, V_grid):
    """
    คำนวณสัมประสิทธิ์การทะลุผ่าน T(E) ด้วย Transfer Matrix Method (TMM)
    """
    N = len(x_grid)
    dx = x_grid[1] - x_grid[0]
    E_joules = E_val * eV
    
    # เวกเตอร์คลื่นในบริเวณทางเข้าและทางออก (V = 0)
    k_in = np.sqrt(2.0 * m_e * E_joules + 0j) / hbar
    k_out = k_in
    
    # เมทริกซ์ส่งผ่านรวมเริ่มต้นเป็นเมทริกซ์เอกลักษณ์ 2x2
    M_total = np.identity(2, dtype=complex)
    
    for i in range(N - 1):
        V_i = V_grid[i] * eV
        k_i = np.sqrt(2.0 * m_e * (E_joules - V_i) + 0j) / hbar
        
        # เมทริกซ์การแพร่ของเฟสภายในชั้นที่ i
        P_i = np.array([[np.exp(-1j * k_i * dx), 0],
                        [0, np.exp(1j * k_i * dx)]], dtype=complex)
        
        # เมทริกซ์รอยต่อระหว่างชั้น i และ i+1
        V_next = V_grid[i+1] * eV
        k_next = np.sqrt(2.0 * m_e * (E_joules - V_next) + 0j) / hbar
        
        d_ij = 0.5 * np.array([[1.0 + k_next/k_i, 1.0 - k_next/k_i],
                              [1.0 - k_next/k_i, 1.0 + k_next/k_i]], dtype=complex)
        
        M_total = M_total @ P_i @ d_ij
        
    # คำนวณสัมประสิทธิ์การทะลุผ่านจากองค์ประกอบ M_total[0, 0]
    T = (np.real(k_out) / np.real(k_in)) * (1.0 / np.abs(M_total[0, 0])**2)
    return np.clip(np.real(T), 0.0, 1.0)

# กำหนดกริดเชิงพื้นที่สำหรับโครงสร้างกำแพงศักย์คู่ RTD (Double Barrier)
L_total = 12.0e-9            # 12 นาโนเมตร
Nx = 1200
x = np.linspace(0, L_total, Nx)
V = np.zeros(Nx)

# สร้างกำแพงศักย์คู่: ความสูง 3.0 eV ความหนา 1.5 nm คั่นด้วยบ่อกว้าง 3.0 nm
V[(x >= 3.0e-9) & (x <= 4.5e-9)] = 3.0
V[(x >= 7.5e-9) & (x <= 9.0e-9)] = 3.0

# สแกนพลังงานของอิเล็กตรอนจาก 0.1 eV ถึง 5.0 eV
energies = np.linspace(0.1, 5.0, 600)
T_spectrum = [transfer_matrix_barrier(E, x, V) for E in energies]

# พล็อตผลลัพธ์
plt.figure(figsize=(10, 6))
plt.plot(energies, T_spectrum, color='#00F0FF', lw=2.2, label='Transmission $T(E)$')
plt.axvline(3.0, color='#FF007F', linestyle='--', alpha=0.7, label='Barrier Height $V_0 = 3.0$ eV')
plt.title('Double-Barrier Resonant Tunneling Diode (RTD): Transmission Spectrum', 
          fontsize=12, fontweight='bold')
plt.xlabel('Incident Electron Energy $E$ (eV)', fontsize=11)
plt.ylabel('Transmission Probability $T(E)$', fontsize=11)
plt.grid(True, linestyle=':', alpha=0.6)
plt.legend(loc='upper left', fontsize=10)
plt.tight_layout()
plt.savefig('rtd_tunneling_simulation.png', dpi=300)
plt.show()
\end{lstlisting}

% ==========================================================
% 4.11 ตัวอย่างการแก้โจทย์ปัญหาเชิงวิเคราะห์แบบ P-S-C-T
% ==========================================================
\section{ตัวอย่างการแก้โจทย์ปัญหาเชิงวิเคราะห์แบบ P-S-C-T}
\label{sec_psct_examples_ch04}

ระเบียบวิธี P-S-C-T เป็นกรอบการคิดวิเคราะห์อย่างเป็นระบบ 4 ขั้นตอน ได้แก่ Problem Statement (ระบุและจำแนกโจทย์), Strategy (วางยุทธศาสตร์และหลักการฟิสิกส์), Calculation (คำนวณและอนุพัทธ์อย่างแม่นยำ), และ Takeaway (ข้อคิดสรุปและบริบทสมัยใหม่)

% ตัวอย่างที่ 1
\begin{psctbox}{การกระจายคลื่นและการฟื้นคืนในบ่อศักย์อนันต์}{ตัวอย่างที่หนึ่ง}
\noindent\textbf{P - Problem Statement (วิเคราะห์โจทย์)} อนุภาคอิเล็กตรอนถูกกักอยู่ในบ่อศักย์อนันต์กว้าง $L = 1.0\text{ nm}$ ณ เวลาเริ่มต้น $t = 0$ ฟังก์ชันคลื่นของระบบอยู่ในสถานะผสมสมมาตร $\Psi(x, 0) = A \left[ \psi_1(x) + \psi_2(x) \right]$ จงหาค่าคงที่นอร์มอลไลเซชัน $A$ คำนวณความน่าจะเป็นที่จะพบอนุภาคในครึ่งซ้ายของบ่อ ($0 \le x \le L/2$) ณ เวลาเริ่มต้น และคำนวณเวลาการสั่นของความน่าจะเป็นรอบหนึ่ง

\vspace{6pt}
\noindent\textbf{S - Strategy \& Principle (วางกลยุทธ์)} ใช้เงื่อนไขความตั้งฉากปรกติของไอเกนสเตต $\braket{\psi_m}{\psi_n} = \delta_{mn}$ เพื่อหาค่า $A$ จากนั้นคำนวณฟังก์ชันความหนาแน่นความน่าจะเป็น $|\Psi(x, t)|^2$ ซึ่งจะมีเทอมการแทรกสอดที่แกว่งตามเวลาด้วยความถี่ $\Delta\omega = (E_2 - E_1)/\hbar$ คาบเวลาการสั่นคือ $T = 2\pi/\Delta\omega$

\vspace{6pt}
\noindent\textbf{C - Calculation \& Proof (การคำนวณและพิสูจน์)}
ขั้นตอนที่หนึ่ง นอร์มอลไลซ์สถานะเริ่มต้น
\begin{equation}
\int_0^L |\Psi(x, 0)|^2 dx = |A|^2 \left[ \int |\psi_1|^2 dx + \int |\psi_2|^2 dx + 2 \operatorname{Re}\int \psi_1^* \psi_2 dx \right] = |A|^2 [1 + 1 + 0] = 2|A|^2 = 1
\end{equation}
จะได้ $A = \frac{1}{\sqrt{2}}$ ดังนั้น $\Psi(x, 0) = \frac{1}{\sqrt{2}} \psi_1(x) + \frac{1}{\sqrt{2}} \psi_2(x)$

ขั้นตอนที่สอง คำนวณความน่าจะเป็นในครึ่งซ้ายของบ่อ
\begin{equation}
P(0 \le x \le L/2) = \int_0^{L/2} |\Psi(x, 0)|^2 dx = \frac{1}{2}\int_0^{L/2} \psi_1^2 dx + \frac{1}{2}\int_0^{L/2} \psi_2^2 dx + \int_0^{L/2} \psi_1(x)\psi_2(x) dx
\end{equation}
เนื่องจาก $\psi_1^2$ และ $\psi_2^2$ มีสมมาตร อินทิกรัลสองเทอมแรกในครึ่งบ่อจะได้ $\frac{1}{2}(0.5) + \frac{1}{2}(0.5) = 0.5$ ส่วนเทอมการแทรกสอด
\begin{align}
\int_0^{L/2} \psi_1 \psi_2 dx &= \frac{2}{L} \int_0^{L/2} \sin\left(\frac{\pi x}{L}\right) \sin\left(\frac{2\pi x}{L}\right) dx \\
&= \frac{1}{L} \int_0^{L/2} \left[ \cos\left(\frac{\pi x}{L}\right) - \cos\left(\frac{3\pi x}{L}\right) \right] dx = \frac{1}{L} \left[ \frac{L}{\pi} (1) - \frac{L}{3\pi} (-1) \right] = \frac{4}{3\pi} \approx 0.424
\end{align}
ดังนั้น ความน่าจะเป็นรวมคือ $P = 0.5 + 0.424 = 0.924$ หรือ $92.4\%$

ขั้นตอนที่สาม คำนวณคาบเวลาการสั่น
\begin{align}
E_1 &= \frac{\pi^2\hbar^2}{2m L^2} = \frac{\pi^2 (1.055 \times 10^{-34})^2}{2(9.109 \times 10^{-31})(1.0 \times 10^{-9})^2} = 6.025 \times 10^{-20}\text{ J} \approx 0.376\text{ eV} \\
E_2 &= 4E_1 = 1.504\text{ eV} \quad \Longrightarrow \quad \Delta E = 3E_1 = 1.808 \times 10^{-19}\text{ J} \\
T_{\text{osc}} &= \frac{2\pi\hbar}{\Delta E} = \frac{6.626 \times 10^{-34}\text{ J}\cdot\text{s}}{1.808 \times 10^{-19}\text{ J}} \approx 3.66 \times 10^{-15}\text{ s} = 3.66\text{ fs}
\end{align}

\vspace{6pt}
\noindent\textbf{T - Takeaway \& Modern Context (ข้อคิดสรุปและบริบทสมัยใหม่)} ความน่าจะเป็นในการพบอนุภาคเอียงไปอยู่ทางครึ่งซ้ายถึง $92.4\%$ เกิดจากการแทรกสอดเชิงบวกของสองสถานะ และกลุ่มคลื่นนี้จะสั่นแกว่งไปมาระหว่างซ้ายและขวาด้วยคาบเวลาเพียง $3.66$ เฟมโตวินาที ความเร็วระดับเฟมโตวินาทีนี้คือมาตรวัดเวลาธรรมชาติของอันตรกิริยาอิเล็กตรอนในโมเลกุล ซึ่งศึกษาได้ด้วยเลเซอร์แอตโตวินาที (Attosecond Physics) ที่ได้รับรางวัลโนเบลปี 2023
\end{psctbox}

% ตัวอย่างที่ 2
\begin{psctbox}{การสะท้อนและการอนุรักษ์กระแสที่รอยต่อศักย์ขั้นบันได}{ตัวอย่างที่สอง}
\noindent\textbf{P - Problem Statement (วิเคราะห์โจทย์)} ลำอนุภาคอิเล็กตรอนพลังงาน $E = 10.0\text{ eV}$ เคลื่อนที่เข้าชนรอยต่อศักย์ขั้นบันไดความสูง $V_0 = 6.0\text{ eV}$ จงคำนวณสัมประสิทธิ์การสะท้อน $R$ สัมประสิทธิ์การทะลุผ่าน $T$ และตรวจสอบกฎการอนุรักษ์กระแสความน่าจะเป็น

\vspace{6pt}
\noindent\textbf{S - Strategy \& Principle (วางกลยุทธ์)} พลังงาน $E > V_0$ ดังนั้นคลื่นสามารถทะลุผ่านเข้าไปในบริเวณที่ II ได้ คำนวณเวกเตอร์คลื่น $k_1$ และ $k_2$ จากนั้นใช้สมการ \eqref{eq_step_refl_coeff} และ \eqref{eq_step_trans_coeff}

\vspace{6pt}
\noindent\textbf{C - Calculation \& Proof (การคำนวณและพิสูจน์)}
ขั้นตอนที่หนึ่ง คำนวณอัตราส่วนเวกเตอร์คลื่น
\begin{equation}
\frac{k_2}{k_1} = \frac{\sqrt{2m(E - V_0)}/\hbar}{\sqrt{2mE}/\hbar} = \sqrt{\frac{E - V_0}{E}} = \sqrt{\frac{10.0 - 6.0}{10.0}} = \sqrt{0.40} \approx 0.63245
\end{equation}

ขั้นตอนที่สอง คำนวณสัมประสิทธิ์การสะท้อน $R$
\begin{equation}
R = \left( \frac{1 - k_2/k_1}{1 + k_2/k_1} \right)^2 = \left( \frac{1 - 0.63245}{1 + 0.63245} \right)^2 = \left( \frac{0.36755}{1.63245} \right)^2 = (0.22515)^2 \approx 0.0507 \quad (5.07\%)
\end{equation}

ขั้นตอนที่สาม คำนวณสัมประสิทธิ์การทะลุผ่าน $T$
\begin{equation}
T = \frac{4(k_2/k_1)}{(1 + k_2/k_1)^2} = \frac{4 \times 0.63245}{(1.63245)^2} = \frac{2.5298}{2.6649} \approx 0.9493 \quad (94.93\%)
\end{equation}

ขั้นตอนที่สี่ ตรวจสอบความสอดคล้อง
\begin{equation}
R + T = 0.0507 + 0.9493 = 1.0000 \quad (100\%)
\end{equation}

\vspace{6pt}
\noindent\textbf{T - Takeaway \& Modern Context (ข้อคิดสรุปและบริบทสมัยใหม่)} แม้อิเล็กตรอนจะมีพลังงานสูงกว่าขั้นบันไดถึง 4 eV แต่ก็ยังมีโอกาสสะท้อนกลับถึง $5.07\%$ ผลลัพธ์นี้ยืนยันว่าการสะท้อนคลื่นควอนตัมไม่ได้ขึ้นอยู่กับสิ่งกีดขวางที่ห้ามไม่ให้อนุภาคผ่าน แต่เกิดจากความไม่ต่อเนื่องของค่าดัชนีหักเหของคลื่นสสาร ซึ่งนำไปใช้เคลือบสารต้านการสะท้อนของอิเล็กตรอนในกล้องจุลทรรศน์อิเล็กตรอน
\end{psctbox}

% ตัวอย่างที่ 3
\begin{psctbox}{การแก้สมการอดิศัยบ่อศักย์จำกัดและระยะแทรกซึม}{ตัวอย่างที่สาม}
\noindent\textbf{P - Problem Statement (วิเคราะห์โจทย์)} อิเล็กตรอนอยู่ในบ่อศักย์จำกัดสมมาตรที่มีความกว้าง $2a = 1.0\text{ nm}$ ($a = 0.5\text{ nm}$) และความลึก $V_0 = 10.0\text{ eV}$ จงคำนวณรัศมีไร้มิติ $R$ หาจำนวนสถานะผูกพันทั้งหมด และคำนวณพลังงานสถานะพื้น $E_1$ พร้อมระยะแทรกซึม $\delta_p$

\vspace{6pt}
\noindent\textbf{S - Strategy \& Principle (วางกลยุทธ์)} คำนวณ $R = \sqrt{2m V_0 a^2/\hbar^2}$ จากนั้นใช้สูตร $N = 1 + \lfloor 2R/\pi \rfloor$ สำหรับสถานะพื้นซึ่งเป็นพาริตีคู่ จะหาคำตอบจากสมการ $\xi\tan\xi = \sqrt{R^2 - \xi^2}$ ด้วยวิธีเชิงตัวเลข

\vspace{6pt}
\noindent\textbf{C - Calculation \& Proof (การคำนวณและพิสูจน์)}
ขั้นตอนที่หนึ่ง คำนวณรัศมีวงกลม $R$
\begin{align}
R &= \frac{a}{\hbar} \sqrt{2m V_0} = \frac{0.5 \times 10^{-9}}{1.055 \times 10^{-34}} \sqrt{2(9.109 \times 10^{-31})(10 \times 1.602 \times 10^{-19})} \\
&= (4.739 \times 10^{24}) \times (1.708 \times 10^{-24}) \approx 8.095
\end{align}

ขั้นตอนที่สอง หาจำนวนสถานะผูกพันทั้งหมด
\begin{equation}
N = 1 + \left\lfloor \frac{2 \times 8.095}{\pi} \right\rfloor = 1 + \lfloor 5.153 \rfloor = 1 + 5 = 6 \text{ สถานะ}
\end{equation}
ประกอบด้วยสถานะคู่ 3 สถานะ และสถานะคี่ 3 สถานะ

ขั้นตอนที่สาม แก้สมการ $\xi\tan\xi = \sqrt{R^2 - \xi^2}$ สำหรับสถานะพื้น
เนื่องจาก $R = 8.095$ ค่อนข้างใหญ่ $\xi$ จะมีค่าเข้าใกล้ $\pi/2 \approx 1.5708$ จากการประมาณค่าหรือทำซ้ำนิวตัน-ราฟสัน จะได้รากของสมการคือ $\xi_1 \approx 1.396\text{ rad}$
\begin{equation}
E_1 + V_0 = \frac{\hbar^2 \xi_1^2}{2m a^2} = \frac{(1.055 \times 10^{-34})^2 (1.396)^2}{2(9.109 \times 10^{-31})(0.5 \times 10^{-9})^2} = 4.76 \times 10^{-20}\text{ J} \approx 0.297\text{ eV}
\end{equation}
ดังนั้น พลังงานผูกพันเทียบกับปากบ่อคือ $E_1 = 0.297 - 10.0 = -9.703\text{ eV}$

ขั้นตอนที่สี่ คำนวณระยะแทรกซึมในบริเวณต้องห้าม
\begin{align}
\kappa_1 &= \frac{\sqrt{2m|E_1|}}{\hbar} = \frac{\sqrt{2(9.109 \times 10^{-31})(9.703 \times 1.602 \times 10^{-19})}}{1.055 \times 10^{-34}} \approx 1.595 \times 10^{10}\text{ m}^{-1} \\
\delta_p &= \frac{1}{\kappa_1} = \frac{1}{1.595 \times 10^{10}} \approx 0.0627\text{ nm} = 0.627\text{ \r{A}}
\end{align}

\vspace{6pt}
\noindent\textbf{T - Takeaway \& Modern Context (ข้อคิดสรุปและบริบทสมัยใหม่)} ฟังก์ชันคลื่นของสถานะพื้นแทรกซึมออกไปนอกขอบบ่อประมาณ $0.063\text{ nm}$ ซึ่งคิดเป็นประมาณ $12\%$ ของความกึ่งกว้างบ่อ ในหลุมควอนตัมของเลเซอร์ GaAs/AlGaAs การแทรกซึมนี้จะกำหนดอันตรกิริยาการซ้อนทับ (Wavefunction Overlap) ระหว่างหลุมควอนตัมข้างเคียงในการสร้างแถบพลังงานย่อย (Superlattice Minibands)
\end{psctbox}

% ตัวอย่างที่ 4
\begin{psctbox}{การกระเจิงและสถานะผูกพันในศักย์เดลตาของดิแรก}{ตัวอย่างที่สี่}
\noindent\textbf{P - Problem Statement (วิเคราะห์โจทย์)} กำหนดบ่อศักย์ดึงดูดเดลตา $V(x) = -\alpha\delta(x)$ โดยที่ $\alpha = 2.0\text{ eV}\cdot\text{nm}$ สำหรับอนุภาคอิเล็กตรอน จงคำนวณพลังงานของสถานะผูกพันเดี่ยว และคำนวณสัมประสิทธิ์การทะลุผ่าน $T$ เมื่อยิงอิเล็กตรอนที่มีพลังงาน $E = 5.0\text{ eV}$ เข้าใส่ศักย์นี้

\vspace{6pt}
\noindent\textbf{S - Strategy \& Principle (วางกลยุทธ์)} พลังงานสถานะผูกพันคำนวณจาก $E_{\text{bound}} = -\frac{m\alpha^2}{2\hbar^2}$ และสัมประสิทธิ์การทะลุผ่านสำหรับคลื่นกระเจิงคำนวณจาก $T = \frac{1}{1 + \frac{m\alpha^2}{2\hbar^2 E}}$

\vspace{6pt}
\noindent\textbf{C - Calculation \& Proof (การคำนวณและพิสูจน์)}
ขั้นตอนที่หนึ่ง คำนวณพารามิเตอร์รวม $\frac{m\alpha^2}{2\hbar^2}$
แปลงค่า $\alpha = 2.0\text{ eV}\cdot\text{nm} = 2.0 \times (1.602 \times 10^{-19}\text{ J}) \times (1.0 \times 10^{-9}\text{ m}) = 3.204 \times 10^{-28}\text{ J}\cdot\text{m}$
\begin{align}
\frac{m\alpha^2}{2\hbar^2} &= \frac{(9.109 \times 10^{-31}\text{ kg})(3.204 \times 10^{-28}\text{ J}\cdot\text{m})^2}{2(1.055 \times 10^{-34}\text{ J}\cdot\text{s})^2} = \frac{9.351 \times 10^{-86}}{2.226 \times 10^{-68}} \\
&= 4.201 \times 10^{-18}\text{ J} = \frac{4.201 \times 10^{-18}}{1.602 \times 10^{-19}} \approx 26.22\text{ eV}
\end{align}

ขั้นตอนที่สอง คำนวณพลังงานสถานะผูกพัน
\begin{equation}
E_{\text{bound}} = -\frac{m\alpha^2}{2\hbar^2} = -26.22\text{ eV}
\end{equation}

ขั้นตอนที่สาม คำนวณสัมประสิทธิ์การทะลุผ่านที่ $E = 5.0\text{ eV}$
\begin{equation}
T = \frac{1}{1 + \frac{m\alpha^2}{2\hbar^2 E}} = \frac{1}{1 + \frac{26.22\text{ eV}}{5.0\text{ eV}}} = \frac{1}{1 + 5.244} = \frac{1}{6.244} \approx 0.160 \quad (16.0\%)
\end{equation}
สัมประสิทธิ์การสะท้อนคือ $R = 1 - T = 84.0\%$

\vspace{6pt}
\noindent\textbf{T - Takeaway \& Modern Context (ข้อคิดสรุปและบริบทสมัยใหม่)} แม้ศักย์เดลตาจะมีความกว้างเป็นศูนย์ในเชิงเรขาคณิต แต่อิเล็กตรอนที่มีพลังงาน 5.0 eV กลับสะท้อนกลับถึง $84\%$ ปรากฏการณ์นี้แสดงให้เห็นว่าความเป็นคลื่นของอนุภาคตอบสนองต่อจุดเอกฐานของสนามไฟฟ้าอย่างรุนแรง ซึ่งเป็นพื้นฐานของการกระเจิงของอิเล็กตรอนโดยอะตอมสิ่งเจือปนเดี่ยว (Single Dopant Scattering) ในสายนาโนซิลิกอน
\end{psctbox}

% ตัวอย่างที่ 5
\begin{psctbox}{การคำนวณกระแสทันเนลลิงในรอยต่อโลหะ-ฉนวน-โลหะ (MIM)}{ตัวอย่างที่ห้า}
\noindent\textbf{P - Problem Statement (วิเคราะห์โจทย์)} รอยต่อรอยต่อนาโนโลหะ-ฉนวน-โลหะ (Metal-Insulator-Metal หรือ MIM) มีชั้นฉนวนออกไซด์หนา $a = 1.2\text{ nm}$ กั้นขวาง มีความสูงกำแพงศักย์ยังผล $V_0 - E_F = 3.2\text{ eV}$ หากความหนาของชั้นฉนวนเพิ่มขึ้น $0.2\text{ nm}$ (กลายเป็น $1.4\text{ nm}$) จงคำนวณว่าความหนาแน่นกระแสทันเนลลิงจะลดลงไปกี่เท่า

\vspace{6pt}
\noindent\textbf{S - Strategy \& Principle (วางกลยุทธ์)} ในขีดจำกัดกำแพงหนา กระแสทันเนลลิงแปรผันตามสัมประสิทธิ์การทะลุผ่าน $J \propto T \approx e^{-2\kappa a}$ โดยที่ $\kappa = \sqrt{2m^*(V_0 - E)}/\hbar$ อัตราส่วนของกระแสคือ $J_2/J_1 = e^{-2\kappa(a_2 - a_1)} = e^{-2\kappa \Delta a}$ กำหนดให้มวลยังผลเท่ากับมวลอิเล็กตรอนอิสระ

\vspace{6pt}
\noindent\textbf{C - Calculation \& Proof (การคำนวณและพิสูจน์)}
ขั้นตอนที่หนึ่ง คำนวณค่าสัมประสิทธิ์การลดทอน $\kappa$
\begin{align}
\kappa &= \frac{\sqrt{2(9.109 \times 10^{-31}\text{ kg})(3.2 \times 1.602 \times 10^{-19}\text{ J})}}{1.055 \times 10^{-34}\text{ J}\cdot\text{s}} = \frac{\sqrt{9.339 \times 10^{-49}}}{1.055 \times 10^{-34}} \\
&= \frac{9.664 \times 10^{-25}}{1.055 \times 10^{-34}} \approx 9.160 \times 10^9\text{ m}^{-1} = 9.160\text{ nm}^{-1}
\end{align}

ขั้นตอนที่สอง คำนวณการเปลี่ยนแปลงเลขชี้กำลัง $2\kappa \Delta a$
\begin{equation}
2\kappa \Delta a = 2 \times (9.160\text{ nm}^{-1}) \times (0.2\text{ nm}) = 3.664
\end{equation}

ขั้นตอนที่สาม คำนวณอัตราส่วนการลดทอนของกระแส
\begin{equation}
\frac{J_2}{J_1} = e^{-2\kappa \Delta a} = e^{-3.664} \approx 0.0256 \approx \frac{1}{39.0}
\end{equation}

\vspace{6pt}
\noindent\textbf{T - Takeaway \& Modern Context (ข้อคิดสรุปและบริบทสมัยใหม่)} การเพิ่มความหนาเพียง $0.2\text{ nm}$ (ขนาดประมาณ 1 ชั้นอะตอม) ส่งผลให้กระแสทันเนลลิงลดฮวบลงถึง 39 เท่า ความไวมหาศาลนี้เป็นหัวใจสำคัญของหน่วยความจำแฟลช (Flash Memory 3D NAND) ซึ่งใช้การกักเก็บประจุผ่านชั้นออกไซด์ และเป็นเหตุผลว่าเหตุใดความหนาแน่นของรอยต่อฉนวนในคิวบิตตัวนำยิ่งยวด (Josephson Junction Transmon Qubit) จึงต้องถูกควบคุมด้วยความแม่นยำระดับซับแองสตรอม
\end{psctbox}

% ตัวอย่างที่ 6
\begin{psctbox}{การคำนวณครึ่งชีวิตการสลายตัวให้อนุภาคแอลฟาตามทฤษฎีกามอฟ}{ตัวอย่างที่หก}
\noindent\textbf{P - Problem Statement (วิเคราะห์โจทย์)} นิวเคลียสพอโลเนียม-212 ($^{212}_{84}\text{Po}$) สลายตัวให้อนุภาคแอลฟาที่มีพลังงานจลน์ $E_\alpha = 8.78\text{ MeV}$ สู่นิวเคลียสตะกั่ว ($^{208}_{82}\text{Pb}$, $Z' = 82$) โดยมีรัศมีนิวเคลียส $R_0 = 1.2 \times (208)^{1/3}\text{ fm} \approx 7.11\text{ fm}$ จงคำนวณจุดวกกลับคลาสสิก $r_c$ เลขชี้กำลังของกามอฟ $2\gamma$ และประมาณค่าครึ่งชีวิต $t_{1/2}$ กำหนดให้ความถี่การชนผนัง $\nu_0 = 1.0 \times 10^{21}\text{ s}^{-1}$

\vspace{6pt}
\noindent\textbf{S - Strategy \& Principle (วางกลยุทธ์)} ใช้ทฤษฎี WKB ของกามอฟ คำนวณ $r_c = \frac{2Z'e^2}{4\pi\varepsilon_0 E_\alpha}$ จากนั้นคำนวณ $2\gamma$ และหา $\lambda = \nu_0 e^{-2\gamma}$ นำไปสู่ $t_{1/2} = \frac{\ln 2}{\lambda}$

\vspace{6pt}
\noindent\textbf{C - Calculation \& Proof (การคำนวณและพิสูจน์)}
ขั้นตอนที่หนึ่ง คำนวณจุดวกกลับคลาสสิก $r_c$
\begin{equation}
r_c = \frac{2 \times 82 \times (1.440\text{ MeV}\cdot\text{fm})}{8.78\text{ MeV}} = \frac{236.16}{8.78} \approx 26.90\text{ fm}
\end{equation}

ขั้นตอนที่สอง คำนวณอัตราส่วน $\sqrt{R_0 / r_c}$
\begin{equation}
\sqrt{\frac{R_0}{r_c}} = \sqrt{\frac{7.11}{26.90}} = \sqrt{0.2643} \approx 0.5141
\end{equation}

ขั้นตอนที่สาม คำนวณความเร็วของอนุภาคแอลฟาและพารามิเตอร์กามอฟ
\begin{align}
v_\alpha &= \sqrt{\frac{2 E_\alpha}{m_\alpha}} = \sqrt{\frac{2 \times (8.78 \times 1.602 \times 10^{-13}\text{ J})}{6.645 \times 10^{-27}\text{ kg}}} \approx 2.057 \times 10^7\text{ m/s} \\
\frac{4 Z' e^2}{4\pi\varepsilon_0 \hbar v_\alpha} &= \frac{4 \times 82 \times (1.602 \times 10^{-19})^2 \times (8.988 \times 10^9)}{(1.055 \times 10^{-34}) \times (2.057 \times 10^7)} \approx 34.80
\end{align}
มุม $\arccos(0.5141) \approx 1.031\text{ rad}$ และ $\sqrt{0.2643(1 - 0.2643)} \approx 0.4410$
\begin{equation}
2\gamma = 34.80 \times [1.031 - 0.4410] = 34.80 \times 0.590 \approx 20.53
\end{equation}

ขั้นตอนที่สี่ คำนวณสัมประสิทธิ์การทะลุผ่านและครึ่งชีวิต
\begin{align}
T &= e^{-20.53} \approx 1.21 \times 10^{-9} \\
\lambda &= \nu_0 T = (1.0 \times 10^{21}\text{ s}^{-1}) \times (1.21 \times 10^{-9}) = 1.21 \times 10^{12}\text{ s}^{-1} \\
t_{1/2} &= \frac{\ln 2}{\lambda} = \frac{0.6931}{1.21 \times 10^{12}\text{ s}^{-1}} \approx 5.7 \times 10^{-7}\text{ s} = 0.57\ \mu\text{s}
\end{align}

\vspace{6pt}
\noindent\textbf{T - Takeaway \& Modern Context (ข้อคิดสรุปและบริบทสมัยใหม่)} ครึ่งชีวิตที่คำนวณได้อยู่ที่ประมาณ $0.57\ \mu\text{s}$ ซึ่งใกล้เคียงกับค่าจริงจากการทดลอง ($0.30\ \mu\text{s}$) อย่างน่าอัศจรรย์ ทฤษฎีนี้ยืนยันว่ากลศาสตร์ควอนตัมทันเนลลิงสามารถก้าวข้ามกำแพงศักย์นิวเคลียร์และอธิบายการแผ่รังสีของสสารได้อย่างสมบูรณ์แบบ
\end{psctbox}

% ตัวอย่างที่ 7
\begin{psctbox}{สเปกโทรสโกปี STS และความหนาแน่นสถานะของแกรฟีน}{ตัวอย่างที่เจ็ด}
\noindent\textbf{P - Problem Statement (วิเคราะห์โจทย์)} ในการวัดสเปกโทรสโกปี STS บนพื้นผิวแกรฟีนชั้นเดี่ยว พบว่าความหนาแน่นสถานะใกล้จุดดิแรกแปรผันตามพลังงานเป็นเชิงเส้น $\rho(E) = \alpha |E - E_D|$ หากแรงดันไบแอส $V_B$ ถูกปรับเปลี่ยนในช่วง $-0.5\text{ V}$ ถึง $+0.5\text{ V}$ จงอธิบายพฤติกรรมของค่าความนำไฟฟ้าเชิงอนุพันธ์ $dI/dV$ และแสดงความสัมพันธ์กับโครงสร้างแถบพลังงาน

\vspace{6pt}
\noindent\textbf{S - Strategy \& Principle (วางกลยุทธ์)} ใช้ทฤษฎีเทอร์ซอฟฟ์-ฮามันน์ ซึ่งระบุว่าค่าความนำไฟฟ้าเชิงอนุพันธ์แปรผันตรงกับความหนาแน่นสถานะเฉพาะที่ $\frac{dI}{dV} \propto \rho_s(E_F + e V_B)$

\vspace{6pt}
\noindent\textbf{C - Calculation \& Proof (การคำนวณและพิสูจน์)}
ขั้นตอนที่หนึ่ง วิเคราะห์ความสัมพันธ์ของกระแสทันเนลลิง
\begin{equation}
I(V_B) \propto \int_0^{e V_B} \rho(E_F + \varepsilon) d\varepsilon = \int_0^{e V_B} \alpha |\varepsilon - \Delta_D| d\varepsilon
\end{equation}
โดยที่ $\Delta_D = E_D - E_F$ คือตำแหน่งของจุดดิแรกเทียบกับระดับแฟร์มี

ขั้นตอนที่สอง หาอนุพันธ์เทียบกับแรงดันไบแอส
\begin{equation}
\frac{dI}{dV_B} \propto e \cdot \rho(E_F + e V_B) = e \alpha |e V_B - \Delta_D|
\end{equation}
กราฟของ $dI/dV_B$ เทียบกับ $V_B$ จะมีรูปร่างเป็นรูปตัวอักษรวี (V-shape) โดยมีจุดต่ำสุดอยู่ที่ $V_{\text{min}} = \Delta_D / e$

\vspace{6pt}
\noindent\textbf{T - Takeaway \& Modern Context (ข้อคิดสรุปและบริบทสมัยใหม่)} เทคนิค STS ช่วยให้นักฟิสิกส์สามารถระบุตำแหน่งของจุดดิแรกและตรวจสอบการมีอยู่ของเฟอร์มิออนไร้มวลแบบดิแรก (Massless Dirac Fermions) ในวัสดุสองมิติได้อย่างแม่นยำ ซึ่งเป็นหัวใจสำคัญของการวิจัยวัสดุควอนตัมแห่งอนาคต
\end{psctbox}

% ==========================================================
% 4.12 สรุปสาระสำคัญประจำบทที่ 4
% ==========================================================
\section{สรุปสาระสำคัญประจำบทที่ 4}
\label{sec_chapter_summary_ch04}

\begin{enumerate}
    \item \textbf{รอยต่อศักย์ขั้นบันได} ---- เมื่ออนุภาคควอนตัมเผชิญกับรอยต่อที่พลังงานศักย์เปลี่ยนแปลงอย่างฉับพลัน จะเกิดการสะท้อนของคลื่นสสารเสมอแม้พลังงานของอนุภาคจะสูงกว่าความสูงของขั้นบันได ($E > V_0$) ซึ่งสอดคล้องกับกฎอนุรักษ์ฟลักซ์กระแสความน่าจะเป็น $R + T = 1$
    
    \item \textbf{คลื่นสลายตัวในบริเวณต้องห้าม} ---- เมื่อ $E < V_0$ ฟังก์ชันคลื่นจะแทรกซึมเข้าไปในเนื้อศักย์แบบเอกซ์โพเนนเชียล $\psi \sim e^{-\kappa x}$ แม้จะเกิดการสะท้อนกลับหมด ($R=1$) แต่ความหนาแน่นความน่าจะเป็นในบริเวณต้องห้ามไม่เป็นศูนย์ ก่อให้เกิดระยะแทรกซึมลักษณะเฉพาะ $\delta_p = 1/(2\kappa)$
    
    \item \textbf{บ่อศักย์อนันต์และพลังงานจุดศูนย์} ---- การกักขังอนุภาคในขอบเขตจำกัดส่งผลให้ระดับพลังงานถูกควอนไตซ์ $E_n = n^2 E_1$ และพลังงานสถานะพื้น $E_1 > 0$ เสมอตามหลักความไม่แน่นอนของไฮเซนแบร์ก
    
    \item \textbf{การฟื้นคืนควอนตัม} ---- เนื่องจากพลังงานแปรผันตาม $n^2$ กลุ่มคลื่นควอนตัมในบ่อศักย์อนันต์จะแผ่กระจายตัวและสามารถกลับมารวมตัวเป็นรูปร่างเดิมได้เมื่อเวลาผ่านไปครบคาบการฟื้นคืน $T_{\text{rev}} = \frac{4mL^2}{\pi\hbar}$
    
    \item \textbf{ความหนาแน่นสถานะในมิติต่างๆ} ---- การลดมิติของระบบควอนตัมเปลี่ยนรูปฟังก์ชันความหนาแน่นสถานะจากพาราโบลาใน 3D สู่ขั้นบันไดใน 2D, ภาวะเอกฐานแวนโฮฟใน 1D, และพีคเดลตาใน 0D (จุดควอนตัม)
    
    \item \textbf{บ่อศักย์จำกัดและสมการอดิศัย} ---- บ่อศักย์จำกัดหนึ่งมิติมีสถานะผูกพันอย่างน้อยหนึ่งสถานะเสมอ (สถานะพาริตีคู่) โดยระดับพลังงานต้องคำนวณจากจุดตัดของสมการอดิศัยกับวงกลมรัศมี $R$
    
    \item \textbf{ศักย์เดลตาของดิแรก} ---- ความไม่ต่อเนื่องของอนุพันธ์อันดับหนึ่ง $\Delta(d\psi/dx) = -\frac{2m\alpha}{\hbar^2}\psi(0)$ นำไปสู่การเกิดสถานะผูกพันเดี่ยว และแบบจำลองบ่อเดลตาคู่ช่วยอธิบายการเกิดพันธะเคมีโคเวเลนต์
    
    \item \textbf{ปรากฏการณ์ควอนตัมทันเนลลิง} ---- อนุภาคสามารถทะลุผ่านกำแพงศักย์ที่สูงกว่าพลังงานของตัวมันเองได้ โดยมีความน่าจะเป็นลดลงแบบเอกซ์โพเนนเชียล $T \propto e^{-2\kappa a}$
    
    \item \textbf{การสั่นพ้องของการทะลุผ่าน} ---- ในกำแพงศักย์คู่หรือเมื่อ $E > V_0$ คลื่นควอนตัมสามารถเกิดการแทรกสอดเสริมกันทำให้กำแพงมีความโปร่งใสสมบูรณ์ ($T=1$) ซึ่งนำไปสู่การทำงานของไดโอด RTD
    
    \item \textbf{การประยุกต์ล้ำยุค} ---- ทฤษฎี WKB ของกามอฟไขปริศนาการสลายตัวให้อนุภาคแอลฟาและปฏิกิริยาฟิวชันในดวงอาทิตย์ ขณะที่ทฤษฎีเทอร์ซอฟฟ์-ฮามันน์เป็นแกนกลางของกล้อง STM และสเปกโทรสโกปี STS ในการควบคุมสสารระดับอะตอม
\end{enumerate}

% ==========================================================
% 4.13 คำถามเชิงมโนทัศน์ประจำบทที่ 4 (20 ข้อ)
% ==========================================================
\section{คำถามเชิงมโนทัศน์ประจำบทที่ 4}
\label{sec_conceptual_questions_ch04}

\begin{enumerate}
    \item จงอธิบายว่าเหตุใดในกลศาสตร์ควอนตัม อนุภาคที่มีพลังงานสูงกว่าขั้นบันไดศักย์ ($E > V_0$) จึงยังมีโอกาสสะท้อนกลับได้ แตกต่างจากฟิสิกส์คลาสสิกอย่างไร
    \item ในการสะท้อนกลับหมดจากขั้นบันไดศักย์เมื่อ $E < V_0$ เหตุใดกระแสความน่าจะเป็นสุทธิในบริเวณต้องห้ามคลาสสิกจึงเป็นศูนย์ แม้ว่าฟังก์ชันคลื่นจะไม่เป็นศูนย์ก็ตาม
    \item ระยะแทรกซึมของฟังก์ชันคลื่นในบริเวณกำแพงศักย์ขึ้นอยู่กับมวลของอนุภาคและความสูงของกำแพงศักย์อย่างไร จงอธิบายเชิงกายภาพ
    \item เหตุใดเลขควอนตัม $n$ ในบ่อศักย์อนันต์จึงไม่สามารถมีค่าเท่ากับศูนย์ ($n \ne 0$) ได้
    \item จงอธิบายความเชื่อมโยงระหว่างพลังงานจุดศูนย์ (Zero-point Energy) ในบ่อศักย์อนันต์กับหลักความไม่แน่นอนของไฮเซนแบร์ก
    \item หากความกว้างของบ่อศักย์อนันต์ขยายตัวเพิ่มขึ้นเป็นสองเท่า ระดับพลังงานของสถานะพื้นและช่องว่างระหว่างระดับพลังงานจะเปลี่ยนแปลงไปอย่างไร
    \item จงอธิบายความแตกต่างระหว่างภาวะเสื่อมเชิงเรขาคณิต (Geometric Degeneracy) และภาวะเสื่อมโดยบังเอิญ (Accidental Degeneracy) ในกล่องศักย์สองมิติ
    \item เหตุใดฟังก์ชันความหนาแน่นสถานะ (Density of States) ในบ่อควอนตัมสองมิติจึงมีลักษณะเป็นขั้นบันไดคงที่ ไม่ขึ้นกับระดับพลังงานภายในแต่ละแถบย่อย
    \item จงพิสูจน์เชิงตรรกะว่าเหตุใดบ่อศักย์จำกัดสมมาตรหนึ่งมิติจึงต้องมีสถานะผูกพันอย่างน้อยหนึ่งสถานะเสมอ ไม่ว่าบ่อจะตื้นหรือแคบเพียงใด
    \item ในบ่อศักย์จำกัด เหตุใดฟังก์ชันคลื่นของสถานะผูกพันจึงมีความยาวคลื่นเดอบรอยล์ภายในบ่อยาวกว่ากรณีบ่อศักย์อนันต์ที่มีความกว้างเท่ากัน
    \item อนุพันธ์อันดับหนึ่งของฟังก์ชันคลื่นมีความต่อเนื่องเสมอหรือไม่ ยกตัวอย่างกรณีที่อนุพันธ์เกิดการกระโดดอย่างฉับพลันพร้อมระบุสาเหตุ
    \item ในแบบจำลองบ่อศักย์เดลตาคู่ เหตุใดระดับพลังงานของสถานะสมมาตรคู่ (Bonding) จึงอยู่ต่ำกว่าระดับพลังงานของสถานะไม่สมมาตร (Antibonding)
    \item ปรากฏการณ์แรมเซาเออร์-ทาวน์เซนด์ (Ramsauer-Townsend Effect) แสดงให้เห็นถึงสมบัติความเป็นคลื่นของอนุภาคอย่างไร
    \item จงอธิบายกลไกการเกิดความต้านทานเชิงอนุพันธ์เป็นลบ (NDR) ในไดโอดทันเนลลิงสั่นพ้อง (RTD)
    \item ในการประมาณค่า WKB เงื่อนไขทางกายภาพใดที่บ่งชี้ว่าศักย์ไฟฟ้ามีความแปรเปลี่ยนช้าเพียงพอที่จะใช้การประมาณนี้ได้
    \item เหตุใดพลังงานของอนุภาคแอลฟาที่เปลี่ยนไปเพียงเล็กน้อย จึงส่งผลให้ครึ่งชีวิตของนิวเคลียสกัมมันตรังสีเปลี่ยนแปลงไปหลายล้านเท่าตามกฎไกเกอร์-นัททอลล์
    \item หากไม่มีปรากฏการณ์ควอนตัมทันเนลลิง ปฏิกิริยานิวเคลียร์ฟิวชันในแกนกลางของดวงอาทิตย์จะสามารถเกิดขึ้นได้หรือไม่ เพราะเหตุใด
    \item กล้องจุลทรรศน์ STM วัดตำแหน่งทางเรขาคณิตของนิวเคลียสอะตอมโดยตรง หรือวัดความหนาแน่นสถานะของอิเล็กตรอน จงอธิบายตามทฤษฎีเทอร์ซอฟฟ์-ฮามันน์
    \item เทคนิคสเปกโทรสโกปี STS อาศัยหลักการใดในการวัดช่องว่างแถบพลังงานของสารกึ่งตัวนำ
    \item เหตุใดอุตสาหกรรมชิปประมวลผลจึงต้องเปลี่ยนผ่านจากทรานซิสเตอร์ FinFET ไปสู่สถาปัตยกรรม RibbonFET GAA ที่ระดับต่ำกว่า 3 นาโนเมตร
\end{enumerate}

% ==========================================================
% 4.14 แบบฝึกหัดคำนวณและพิสูจน์เชิงวิเคราะห์ประจำบทที่ 4 (20 ข้อ)
% ==========================================================
\section{แบบฝึกหัดคำนวณและพิสูจน์เชิงวิเคราะห์ประจำบทที่ 4}
\label{sec_analytical_problems_ch04}

\begin{enumerate}
    \item อนุภาคอิเล็กตรอนพลังงาน $E = 8.0\text{ eV}$ เคลื่อนที่เข้าชนรอยต่อศักย์ขั้นบันได $V_0 = 5.0\text{ eV}$ จงคำนวณหาแอมพลิจูดการสะท้อน $B/A$ สัมประสิทธิ์การสะท้อน $R$ และสัมประสิทธิ์การทะลุผ่าน $T$
    
    \item สำหรับรอยต่อศักย์ขั้นบันไดที่มี $E < V_0$ จงพิสูจน์ว่าเฟสชิฟต์ $\phi$ ของคลื่นสะท้อนสอดคล้องกับความสัมพันธ์ $\cos(2\phi) = \frac{2E - V_0}{V_0}$
    
    \item อิเล็กตรอนถูกกักอยู่ในบ่อศักย์อนันต์ความกว้าง $L = 0.20\text{ nm}$ จงคำนวณพลังงานสถานะพื้น $E_1$ และความยาวคลื่นของโฟตอนที่ปล่อยออกมาเมื่ออิเล็กตรอนเปลี่ยนระดับพลังงานจาก $n=3$ ลงสู่ $n=1$
    
    \item จงพิสูจน์ว่าสำหรับอนุภาคในบ่อศักย์อนันต์หนึ่งมิติ ค่าคาดหวังของโมเมนตัมกำลังสี่มีค่าเท่ากับ $\langle p^4 \rangle = \left(\frac{n\pi\hbar}{L}\right)^4$
    
    \item จงคำนวณความน่าจะเป็นที่จะพบอนุภาคในสถานะ $n=1$ ของบ่อศักย์อนันต์ในช่วงระหว่าง $x = 0.49L$ ถึง $x = 0.51L$ และเปรียบเทียบกับผลการทำนายทางคลาสสิก
    
    \item ในกล่องศักย์สี่เหลี่ยมจัตุรัสสองมิติกว้างด้านละ $L$ จงหาระดับพลังงานและความเสื่อม (Degeneracy) ของระดับพลังงาน 5 ระดับแรก
    
    \item พิจารณากล่องศักย์สามมิติรูปลูกบาศก์กว้าง $L$ จงหาระดับพลังงานที่มีภาวะเสื่อม 6 เท่า (6-fold Degenerate) ที่มีค่าพลังงานต่ำที่สุด พร้อมระบุชุดเลขควอนตัม $(n_x, n_y, n_z)$ ทั้งหมด
    
    \item อนุภาคอยู่ในบ่อศักย์จำกัดสมมาตรที่มี $V_0 a^2 = \frac{9\pi^2\hbar^2}{8m}$ จงหาจำนวนสถานะผูกพันทั้งหมด พร้อมระบุว่าเป็นสถานะคู่กี่สถานะและสถานะคี่กี่สถานะ
    
    \item สำหรับบ่อศักย์จำกัดตื้นมาก ($R \ll 1$) จงพิสูจน์ว่าพลังงานของสถานะผูกพันเดี่ยวมีค่าประมาณ $E \approx -V_0 \left(\frac{m V_0 a^2}{2\hbar^2}\right)$
    
    \item อนุภาคอิเล็กตรอนเผชิญกำแพงศักย์สี่เหลี่ยมกว้าง $a = 0.40\text{ nm}$ และสูง $V_0 = 12.0\text{ eV}$ จงคำนวณสัมประสิทธิ์การทะลุผ่าน $T$ เมื่ออิเล็กตรอนมีพลังงาน $E = 2.0\text{ eV}$
    
    \item จงพิสูจน์ว่าในขีดจำกัดที่กำแพงศักย์สี่เหลี่ยมบางมากแต่สูงมาก โดยที่ $V_0 a \to \alpha$ สัมประสิทธิ์การทะลุผ่านของกำแพงสี่เหลี่ยมจะลดรูปลงสู่สูตรการทะลุผ่านของศักย์เดลตาของดิแรก
    
    \item อนุภาคตกกระทบกำแพงศักย์เดลตาผลัก $V(x) = +\alpha\delta(x)$ ด้วยพลังงาน $E$ จงหาพลังงาน $E$ ที่ทำให้สัมประสิทธิ์การสะท้อนมีค่าเท่ากับสัมประสิทธิ์การทะลุผ่านพอดี ($R = T = 0.5$)
    
    \item ในแบบจำลองบ่อศักย์เดลตาคู่ $V(x) = -\alpha[\delta(x+a) + \delta(x-a)]$ จงแสดงว่าเมื่อระยะห่าง $2a \to \infty$ ระดับพลังงานของสถานะพันธะและปฏิพันธะจะลู่เข้าสู่ค่าพลังงานของบ่อเดลตาเดี่ยว
    
    \item ลำอิเล็กตรอนพลังงาน $E = 54\text{ eV}$ พุ่งชนกำแพงศักย์คู่ที่มีบ่อคั่นกลาง จงคำนวณความกว้างของบ่อที่ทำให้เกิดการสั่นพ้องของการทะลุผ่านสมบูรณ์แบบ ($T=1$) สำหรับโหมด $n=2$
    
    \item จงใช้การประมาณค่า WKB คำนวณสัมประสิทธิ์การทะลุผ่านของกำแพงศักย์รูปสามเหลี่ยม $V(x) = V_0 (1 - x/a)$ สำหรับ $0 \le x \le a$
    
    \item จงพิสูจน์เงื่อนไขการควอนไตซ์แบบบอร์-ซอมเมอร์เฟลด์ $\oint p dx = (n + 1/2)2\pi\hbar$ จากการเชื่อมต่อฟังก์ชันคลื่น WKB ข้ามจุดวกกลับคลาสสิกสองจุด
    
    \item นิวเคลียส $^{238}_{92}\text{U}$ ปล่อยอนุภาคแอลฟาพลังงาน $4.20\text{ MeV}$ สู่นิวเคลียสทอเรียม ($Z'=90$) หากรัศมีนิวเคลียสคือ $R_0 = 7.5\text{ fm}$ จงคำนวณเลขชี้กำลังกามอฟ $2\gamma$ และประมาณค่าครึ่งชีวิตในหน่วยปี
    
    \item ปลายเข็ม STM ทำจากทังสเตนมีฟังก์ชันงาน $\Phi = 4.5\text{ eV}$ อยู่เหนือพื้นผิวทองคำ หากต้องการให้กระแสทันเนลลิงลดลงเหลือ $1\%$ ของค่าเดิม ปลายเข็มจะต้องถูกถอยร่นออกไปเป็นระยะทางเท่าใด
    
    \item ในการวัด STS บนตัวนำยิ่งยวด พบว่าความหนาแน่นสถานะมีรูปแบบตามทฤษฎี BCS $\rho(E) = \rho_0 \frac{|E|}{\sqrt{E^2 - \Delta^2}}$ สำหรับ $|E| > \Delta$ จงร่างกราฟ $dI/dV$ เทียบกับแรงดันไบแอส $V$
    
    \item พิจารณากลุ่มคลื่นเกาส์เซียนที่เคลื่อนที่ชนกำแพงศักย์ จงพิสูจน์โดยใช้ทฤษฎีบทการอนุรักษ์ความน่าจะเป็นว่า อัตราการเปลี่ยนแปลงของความน่าจะเป็นรวมภายในบริเวณกำแพงศักย์มีค่าเท่ากับผลต่างของฟลักซ์กระแสที่ไหลเข้าและไหลออก $\frac{d}{dt}\int_0^a |\psi|^2 dx = J(0, t) - J(a, t)$
\end{enumerate}
